\exercisesfor{2}

\begin{enumerate}
\def\labelenumi{\arabic{enumi}.}
\tightlist
\item
  Let g be a free Lie algebra with basic family \((x_{1},\ldots,x_{n},\ldots)\) . Let \(g_{n}\) denote the subalgebra of g generated by \(x_{1},\ldots,x_{n}\) and let \(b_{n}\) denote the smallest ideal of \(g_{n}\) containing \(x_{n}\) .
\end{enumerate}

\begin{enumerate}
\def\labelenumi{(\alph{enumi})}
\tightlist
\item
  Show that \(\mathfrak{b}_n\) is the submodule of \(\mathfrak{g}_n\) generated by the
\end{enumerate}

\[
\left(\operatorname{ad} \left(x _ {i _ {1}}\right) \circ \dots \circ \operatorname{ad} \left(x _ {i _ {k}}\right)\right) \left(x _ {n}\right),
\]

where \(k \geqslant 0\) and \(i_h \leqslant n\) for all \(h\) .

\begin{enumerate}
\def\labelenumi{(\alph{enumi})}
\setcounter{enumi}{1}
\item
  Show that \(\mathfrak{g}_n = \mathfrak{g}_{n-1} \oplus \mathfrak{b}_n\) ; deduce that \(\mathfrak{g}\) is the direct sum of the \(\mathfrak{b}_n\) for \(n \geqslant 1\) .
\item
  Show, using (a) and (b), that the K-module g is generated by the elements \([x_{i_{1}}, [x_{i_{2}}, \ldots, [x_{i_{k-1}}, x_{i_{k}}] \ldots]]\) , where \(k \geqslant 0\) and \(i_{h} \leqslant i_{k}\) for h \textless{} k.
\end{enumerate}

¶ 2. Let X be a countable set with at least two elements and let H be the set of subsets of M(X) which are Hall sets (cf.~Definition 2). Show that Card(H) = \(2^{N_0}\).

\begin{enumerate}
\def\labelenumi{\arabic{enumi}.}
\setcounter{enumi}{2}
\item
  Let \(\mathbf{X}\) be a set with a total ordering. Show that there exists a Hall set relative to \(\mathbf{X}\) such that \(\mathrm{H} \cap \mathrm{M}^3(\mathbf{X})\) consists of the products \(z(yx)\) with \(y < x\) , \(y \leqslant z\) and that \(\mathrm{H} \cap \mathrm{M}^4(\mathbf{X})\) consists of the products \(w(z(yx))\) with \(w \geqslant z \geqslant y\) , \(y < x\) and the products \((ab)(cd)\) with \(a < b\) , \(c < d\) and either \(a < c\) or \(a = c\) and \(b < d\) .
\item
  \begin{enumerate}
  \def\labelenumii{(\alph{enumii})}
  \tightlist
  \item
    Show that the Lie algebra defined by the presentation
  \end{enumerate}
\end{enumerate}

\[
\{x, y; [ x, [ x, y ] ] = [ y, [ x, y ] ] = 0 \}
\]

admits the basis \((x,y,[x,y])\) . Define an embedding of this algebra in a matrix algebra.

\begin{enumerate}
\def\labelenumi{(\alph{enumi})}
\setcounter{enumi}{1}
\tightlist
\item
  Do the same for the Lie algebra defined by the presentation
\end{enumerate}

\[
\{x, y; [ x, [ x, y ] ] - 2 x = [ y, [ x, y ] ] + 2 y = 0 \}.
\]

\begin{enumerate}
\def\labelenumi{(\alph{enumi})}
\setcounter{enumi}{2}
\tightlist
\item
  Show that the Lie algebra with presentation
\end{enumerate}

\[
\{x, y, z; [ x, y ] - x = [ y, z ] - y = [ z, x ] - z = 0 \}
\]

reduces to 0.

\begin{enumerate}
\def\labelenumi{\arabic{enumi}.}
\setcounter{enumi}{4}
\item
  Let \(\mathfrak{g}\) be a free Lie algebra and let \(\mathfrak{r}\) be an ideal of \(\mathfrak{g}\) . Suppose that \(\mathfrak{r} = [\mathfrak{g},\mathfrak{r}]\) . Show that \(\mathfrak{r} = \{0\}\) (use the fact that \(\bigcap_{n\geqslant 0}\mathcal{C}^n\mathfrak{g} = \{0\}\) ).
\item
  Let \(\mathbf{E}\) be a module. Let \(\mathbf{ME}\) be the free algebra of \(\mathbf{E}\) (Algebra, Chapter III, §2, Exercise 13). Recall that \(\mathrm{ME} = \bigoplus_{n\geqslant 1}\mathrm{E}_n\) , where
\end{enumerate}

\[
\mathrm{E} _ {1} = \mathrm{E}, \qquad \mathrm{E} _ {n} = \bigoplus_ {p + q = n} \mathrm{E} _ {p} \otimes \mathrm{E} _ {q} \quad \text {if} \quad n \geqslant 2,
\]

and the algebra structure on ME is defined by means of the canonical mappings \(\mathbf{E}_n\otimes \mathbf{E}_m\to \mathbf{E}_{n + m}\) .

\begin{enumerate}
\def\labelenumi{(\alph{enumi})}
\item
  Let J be the smallest two-sided ideal of ME containing the elements of the form xx and \((xy)z + (yz)x + (zx)y\) with x, y, z in E. Let LE = ME/J. Show that J is a graded ideal of ME; deduce a graduation \((\mathrm{L}^{n}\mathrm{E})_{n\geqslant1}\) of LE.
\item
  Show that LE is a Lie algebra; it is called the free Lie algebra of the module E. Show that for every Lie algebra g and every linear mapping f: \(E \to g\) there exists one and only one Lie algebra homomorphism F: \(LE \to g\) which extends f.
\item
  Define isomorphisms \(\mathbf{E} \to \mathbf{L}^1\mathbf{E}\) and \(\wedge^2\mathbf{E} \to \mathbf{L}^2\mathbf{E}\) . Show that \(\mathbf{L}^3\mathbf{E}\) is identified with the quotient of \(\mathbf{E} \otimes \wedge^2\mathbf{E}\) by the submodule generated by the elements
\end{enumerate}

\[
x \otimes (y \wedge z) + y \otimes (z \wedge x) + z \otimes (x \wedge y) \quad \text { for } x, y, z \text { in } E.
\]

\begin{enumerate}
\def\labelenumi{(\alph{enumi})}
\setcounter{enumi}{3}
\item
  Show that, when E is a free module of basis X, LE is identified with the free Lie algebra \(\mathrm{L}(\mathbf{X})\) defined in no. 2.
\item
  Show that the canonical mapping of E into the enveloping algebra U(LE) of LE can be extended to an isomorphism of the tensor algebra TE onto U(LE).
\item
  Let \(\sigma\) be the linear mapping of \(\wedge^2\mathrm{E}\) into \(\mathsf{T}^2\mathsf{E}\) such that
\end{enumerate}

\[
\sigma (x \wedge y) = x \otimes y - y \otimes x.
\]

Construct a module E for which \(\sigma\) is not injective. Deduce that for this module the canonical mapping of LE into U(LE) is not injective (compare with Exercise 9 of Chapter I, §2).

\begin{enumerate}
\def\labelenumi{\arabic{enumi}.}
\setcounter{enumi}{6}
\tightlist
\item
  Suppose that K is a field. Let l be a free Lie algebra with basic family \((x_{i})_{i\in I}\) and let M be an l-module.
\end{enumerate}

\begin{enumerate}
\def\labelenumi{(\alph{enumi})}
\item
  Show that \(\mathbf{H}^2 (\mathbf{l},\mathbf{M}) = \{0\}\) , cf.~Chapter I, § 3, Exercise 12. (Use Corollary 2 to Proposition 1 and part (i) of the Exercise in question.)
\item
  Show that, for every family \((m_{i})_{i\in I}\) of elements of M, there exists a cocycle \(\phi:l\to M\) of degree 1 such that \(\phi (x_i) = m_i\) and that this cocycle is unique. Deduce an exact sequence:
\end{enumerate}

\[
0 \rightarrow \mathrm{H} ^ {0} (\mathfrak {l}, \mathrm{M}) \rightarrow \mathrm{M} \rightarrow \mathrm{M} ^ {\mathrm{I}} \rightarrow \mathrm{H} ^ {1} (\mathfrak {l}, \mathrm{M}) \rightarrow 0.
\]

If I is finite of cardinal \(n\) and M is of finite rank over K, then

\[
\operatorname{rg} \mathrm{H} ^ {1} (\mathfrak {l}, \mathrm{M}) - \operatorname{rg} \mathrm{H} ^ {0} (\mathfrak {l}, \mathrm{M}) = (n - 1) \operatorname{rg} \mathrm{M}.
\]

\begin{enumerate}
\def\labelenumi{\arabic{enumi}.}
\setcounter{enumi}{7}
\tightlist
\item
  Suppose that K is a field. Let l be a free Lie algebra with basic family \((\overline{x}_{i})_{i\in I}\) , let r be an ideal of l contained in {[}l, l{]} and let g = l/r; let \(x_{i}\) denote the image of \(\overline{x}_{i}\) in g.
\end{enumerate}

Show the equivalence of the following properties:

\begin{enumerate}
\def\labelenumi{(\roman{enumi})}
\item
  g is a free Lie algebra.
\item
  g admits \((x_{i})\) as basic family (i.e.~r = \{0\}).
\item
  For every g-module M, H²(g, M) = \{0\}.
\item
  If g operates trivially on K, then \(H^{2}(g, K) = \{0\}\) .
\end{enumerate}

(The implications (ii) \(\Rightarrow\) (i) and (iii) \(\Rightarrow\) (iv) are obvious and (i) \(\Rightarrow\) (iii) follows from Exercise 7. To prove that (iv) \(\Rightarrow\) (ii) it suffices to show that \(\mathfrak{r} = [\mathfrak{l},\mathfrak{r}]\) , cf.~Exercise 5; otherwise take a hyperplane \(\mathfrak{h}\) of \(\mathfrak{r}\) containing \([\mathfrak{l},\mathfrak{r}]\) and note that the extension \(\mathfrak{r} / \mathfrak{h}\to \mathfrak{l} / \mathfrak{h}\to \mathfrak{l} / \mathfrak{r} = \mathfrak{g}\) is essential.)

¶ 9. Let \(\mathfrak{g} = \bigoplus_{n \geqslant 1} \mathfrak{g}_n\) be a graded Lie algebra. If \(\mathfrak{h}\) is a graded subalgebra of \(\mathfrak{g}\) such that \(\mathfrak{g} = \mathfrak{h} + [\mathfrak{g}, \mathfrak{g}]\) , show that \(\mathfrak{h} = \mathfrak{g}\) . If \((x_i)\) is a family of homogeneous elements of \(\mathfrak{g}\) , deduce that \((x_i)\) is a generating family if and only if the images of the \(x_i\) in \(\mathfrak{g}/[\mathfrak{g}, \mathfrak{g}]\) generate \(\mathfrak{g}/[\mathfrak{g}, \mathfrak{g}]\) .

If K is a field on which g operates trivially and \(\mathrm{H}^{2}(\mathfrak{g},\mathbf{K})=\{0\}\) , show that the family \((x_{i})\) is basic if and only if the images of the \(x_{i}\) in g/{[}g, g{]} form a basis of this vector space. (Use Exercise 8.)†

\begin{enumerate}
\def\labelenumi{\arabic{enumi}.}
\setcounter{enumi}{9}
\tightlist
\item
  Let \(t\) be a free Lie algebra admitting a finite basic family with \(n\) elements and let \(u\) be a surjective endomorphism of \(l\) . Show that \(u\) is an isomorphism. (Write \(\mathrm{gr}^n(l) = \mathcal{C}^n / \mathcal{C}^{n+1}l\) and denote by \(\mathrm{gr}^n(u)\) the endomorphism of \(\mathrm{gr}^n(l)\) derived from \(u\) ; note that \(\mathrm{gr}^n(u)\) is surjective; as \(\mathrm{gr}^n(l)\) is a free K-module of finite rank, deduce that \(\mathrm{gr}^n(u)\) is bijective since the kernel of \(u\) is contained in \(\bigcap_{n} \mathcal{C}^n l\) , which is \(\{0\}\) .)
\end{enumerate}

If \((y_{1},\ldots ,y_{m})\) is a generating family of l, show that \(m\geqslant n\) and that we have equality if and only if \((y_{1},\ldots ,y_{m})\) is a basic family.

¶ 11. Let Y and Z be two disjoint sets (with Y totally ordered), let \(X = Y \cup Z\) and \(l = L(X)\) . Let \(a_{Y}\) be the submodule of l generated by Z and {[}l, l{]}; it is an ideal of the Lie algebra l. We propose to find explicitly a basic family of \(a_{Y}\) .

Let \(\mathbf{M}\) be the set of positive integral functions on \(\mathbf{Y}\) of finite support. If \(m\in \mathbf{M}\) , we denote by \(\theta^m\) the endomorphism of l given by

\[
\theta^ {m} = \prod_ {y \in Y} (a d y) ^ {m (y)},
\]

where the product is relative to the order relation given on Y. Let \(\mathfrak{S}\) denote the subset of \(\mathbf{Y} \times \mathbf{M}\) consisting of the ordered pairs \((u, m)\) such that there exists \(y \in \mathbf{Y}\) for which \(y > u\) and \(m(y) \geqslant 1\) . Show that the elements

\[
\theta^ {m} (u) \quad \text { for } \quad (u, m) \in \mathfrak {S} \quad \text { and } \quad \theta^ {m} (z) \quad \text { for } \quad (z, m) \in Z \times M
\]

form a basic family of the Lie algebra \(\mathfrak{a}_{\mathbf{Y}}\) .

(Reduce the problem to the case where Y is finite and argue by induction on Card(Y). Use the Corollary to Proposition 10 applied to the greatest element y of Y and apply the induction hypothesis to \(Y = \{y\}\) .)

\begin{enumerate}
\def\labelenumi{\arabic{enumi}.}
\setcounter{enumi}{11}
\tightlist
\item
  Let \(\mathbf{X} = \{x, y\}\) be a set of two elements. Show that the derived algebra of \(\mathbf{L}(\mathbf{X})\) is a free Lie algebra admitting as basic family the elements \(((\operatorname{ad}y)^q \circ (\operatorname{ad}x)^p)(y)\) , for \(p \geqslant 1\) , \(q \geqslant 0\) . (Use Exercise 11.)
\end{enumerate}

¶ 13. (In this exercise we assume that every projective module over K is free; this is so for example if K is a principal ideal domain.)

Let \(\mathfrak{l} = \sum_{n=1}^{\infty} \mathfrak{l}_n\) be a graded Lie algebra admitting a basic family B consisting of homogeneous elements and let \(\mathfrak{h}\) be a graded Lie subalgebra of \(\mathfrak{l}\) , which is a direct factor of \(\mathfrak{l}\) considered as a module.

\begin{enumerate}
\def\labelenumi{(\alph{enumi})}
\tightlist
\item
  For all \(i \geqslant 0\) , let \(\mathfrak{l}^{(i)}\) be the graded subalgebra of \(\mathfrak{l}\) such that
\end{enumerate}

\[
\begin{array}{l l} \mathfrak {l} _ {j} ^ {(i)} = \mathfrak {h} _ {j} & \text { if } j \leqslant i. \\ \mathfrak {l} _ {j} ^ {(i)} = \mathfrak {l} _ {j} & \text { if } j > i. \end{array}
\]

Then \(l = \ell^{(0)} \supset \ell^{(1)} \supset \cdots \supset \mathfrak{h}\) . Show, arguing by induction on \(i\) , the existence of a basic family \(B^{(i)}\) of \(\ell^{(i)}\) , consisting of homogeneous elements, such that the elements of \(B^{(i-1)}\) and of \(B^{(i)}\) of degree \(< i - 1\) are the same. (Suppose that \(B^{(i-1)}\) has been constructed. Let \(m_i\) be the intersection of \(l_i = l_i^{(i-2)}\) with the subalgebra generated by the \(b_j\) for \(j < i\) and let \(b_i\) be the submodule of \(l_i\) generated by the elements of \(B^{(i-1)}\) of degree \(i\) . The induction hypothesis implies that \(l_i = m_i \oplus b_i\) . As \(m_i \subset b_i\) , we can decompose \(b_i\) as \(b_i = y_i \oplus b_i\) , so that \(b_i = m_i \oplus b_i\) ; by the hypothesis on K the modules \(y_i\) and \(b_i\) are free;

changing the elements of \(\mathrm{B}^{(i-1)}\) of degree i if necessary, we can assume that \(b_{i}\) is generated by a subset of \(\mathrm{B}^{(i-1)}\) . Applying Exercise 11 to \(\mathfrak{l}^{(i-1)}\) and its ideal \(\mathfrak{l}^{(i)}\) , we obtain a basic family \(\mathrm{B}^{(i)}\) of \(\mathfrak{l}^{(i)}\) with the desired properties.)

\begin{enumerate}
\def\labelenumi{(\alph{enumi})}
\setcounter{enumi}{1}
\tightlist
\item
  Deduce from (a) the fact that \(\mathfrak{h}\) has a basic family consisting of homogeneous elements.
\end{enumerate}

¶ 14. Suppose that K is a field. Let X be a set and x, y two elements of L(X) which are linearly independent over K. Show that the family \((x,y)\) is free in the Lie algebra L(X). (Let \(x_{p}\) (resp. \(y_{q}\) ) be the non-zero homogeneous component of x (resp. y) of highest degree. Adding if necessary a multiple of y to x, we can assume that \(x_{p}\) and \(y_{q}\) are linearly independent. The subalgebra of L(X) generated by \(x_{p}\) and \(y_{q}\) is graded and hence free, cf.~Exercise 13. Deduce that \((x_{p},y_{q})\) is a free family and pass from this to \((x,y)\) .)

¶ 15. Let \(\mathrm{X} = \{x, y\}\) be a set with two elements and let \(\sigma\) be an automorphism of the Lie algebra \(\mathrm{L}(\mathrm{X})\) . Show that, if \(\mathrm{K}\) is a field, \(\sigma\) preserves the graduation of \(\mathrm{L}(\mathrm{X})\) (apply Exercise 14 to the non-zero homogeneous components of \(\sigma(x)\) and \(\sigma(y)\) of highest degree); deduce an isomorphism of \(\operatorname{Aut}(\mathrm{L}(\mathrm{X}))\) onto \(\mathbf{GL}(2, \mathrm{K})\) . Extend these results to the case where the ring \(\mathrm{K}\) has no nilpotent element \(\neq 0\) . When \(\mathrm{K}\) contains an element \(\varepsilon \neq 0\) of zero square, show that \(x \mapsto x, y \mapsto y + \varepsilon[x, y]\) can be extended to an automorphism of \(\mathrm{L}(\mathrm{X})\) which does not preserve the graduation of \(\mathrm{L}(\mathrm{X})\) .

¶ 16. Suppose that K is a Q-algebra. If g is a Lie algebra, we denote by U(g) its enveloping algebra, by σ the canonical mapping of g into U(g) and by S(g) the symmetric algebra of the K-module g. There exists one and only one linear mapping

\[
\eta (\mathfrak {g}): \mathrm{S} (\mathfrak {g}) \rightarrow \mathrm{U} (\mathfrak {g})
\]

such that \(\eta (\mathfrak{g})(x^n) = \sigma (x)^n\) for all \(x\in \mathfrak{g}\) and all \(n\in \mathbf{N}\) . Then

\[
\eta (\mathfrak {g}) (x _ {1} \dots x _ {n}) = \frac {1}{n !} \sum_ {s \in \mathfrak {S} _ {n}} \sigma (x _ {s (1)} \dots x _ {s (n)}), \quad x _ {i} \in \mathfrak {g}.
\]

We propose to show that \(\eta(g)\) is bijective.

\begin{enumerate}
\def\labelenumi{(\alph{enumi})}
\item
  Prove that \(\eta(g)\) is surjective and that it is bijective when the K-module \(g\) is free (use the Poincaré-Birkhoff-Witt Theorem).
\item
  Let \(\mathfrak{h}\) be an ideal of \(\mathfrak{g}\) . Let \(\mathfrak{s}_{\mathfrak{h}}\) (resp. \(\mathfrak{u}_{\mathfrak{h}}\) ) denote the ideal of \(S(\mathfrak{g})\) (resp. the two-sided ideal of \(U(\mathfrak{g})\) ) generated by \(\mathfrak{h}\) (resp. \(\sigma(\mathfrak{h})\) ). The diagram
\end{enumerate}

\[
\begin{array}{c} 0 \to s _ {\mathfrak {h}} \to S (g) \to S (g / \mathfrak {h}) \to 0 \\ \downarrow^ {\eta (g)} \quad \downarrow^ {\eta (g / \mathfrak {h})} \\ 0 \to u _ {\mathfrak {h}} \to U (g) \to U (g / \mathfrak {h}) \to 0 \end{array}
\]

is commutative and its rows are exact. Deduce that the image \(\tilde{u}_{\mathfrak{h}}\) of \(s_{\mathfrak{h}}\) under \(\eta(g)\) is contained in \(u_{\mathfrak{h}}\) and that \(\tilde{u}_{\mathfrak{h}} = u_{\mathfrak{h}}\) when \(\eta(g)\) and \(\eta(g/h)\) are injective (in particular, when the module \(g\) and \(g/h\) are free).

\begin{enumerate}
\def\labelenumi{(\alph{enumi})}
\setcounter{enumi}{2}
\item
  Take \(\mathfrak{g}\) to be a free Lie algebra with basic family \((\mathbf{X}_1, \ldots, \mathbf{X}_n, \mathrm{H})\) , where \(n \geqslant 0\) , and take \(\mathfrak{h}\) to be the ideal of \(\mathfrak{g}\) generated by H. Show that \(\mathfrak{g}\) and \(\mathfrak{g} / \mathfrak{h}\) are free modules. Deduce that, in that case, \(\tilde{\mathfrak{u}}_{\mathfrak{h}} = \mathfrak{u}_{\mathfrak{h}}\) . In particular, there exists \(x \in \mathfrak{s}_{\mathfrak{h}}\) such that \((\eta(\mathfrak{g}))(x) = \sigma(\mathbf{X}_1) \ldots \sigma(\mathbf{X}_n) \sigma(\mathrm{H})\) .
\item
  We return to the general case. Show that \(\mathfrak{u}_{\mathfrak{h}}\) is generated over \(\mathrm{K}\) by the elements of the form \(\sigma(x_1) \ldots \sigma(x_n)\sigma(h)\) , with \(n \geqslant 0\) , \(x_i \in \mathfrak{g}\) and \(h \in \mathfrak{h}\) (note that \(\mathfrak{u}_{\mathfrak{h}}\) coincides with the left ideal generated by \(\mathfrak{h}\) ). Show, using (c), that such an element belongs to \(\tilde{\mathfrak{u}}_{\mathfrak{h}}\) (use a suitable homomorphism of a free Lie algebra into \(\mathfrak{g}\) ). Deduce that \(\tilde{\mathfrak{u}}_{\mathfrak{h}} = \mathfrak{u}_{\mathfrak{h}}\) .
\item
  Show that, if \(\eta(g)\) is bijective, so is \(\eta(g/h)\) . Deduce finally that, for every Lie algebra \(g, \eta(g)\) is bijective (write \(g\) as the quotient of a free Lie algebra) and that the canonical homomorphism
\end{enumerate}

\[
\omega : \mathsf {S} (\mathfrak {g}) \to \operatorname{gr} \mathrm{U} (\mathfrak {g}) \quad (\text { cf.   Chapter   I,   } \S 2, \text {   no.   } 6)
\]

is an isomorphism (``Poincaré-Birkhoff-Witt Theorem'' for Lie algebras over Q-algebras).

\exercisesfor{3}

The letter X denotes a set.

\begin{enumerate}
\def\labelenumi{\arabic{enumi}.}
\tightlist
\item
  \begin{enumerate}
  \def\labelenumii{(\alph{enumii})}
  \tightlist
  \item
    Show that every element \(u \in \mathrm{A}^{+}(\mathrm{X})\) can be written uniquely in the form \(u = \sum_{x \in \mathrm{X}} u_{x} x\) , where \(u_{x} \in \mathrm{A}(\mathrm{X})\) .
  \end{enumerate}
\end{enumerate}

\begin{enumerate}
\def\labelenumi{(\alph{enumi})}
\setcounter{enumi}{1}
\tightlist
\item
  Show that \(\{0\}\) is the only submodule of \(\mathrm{A}^{+}(\mathrm{X})\) which is stable under all the mappings \(u \mapsto u_{x} (x \in \mathrm{X})\) . (If a is such a submodule and \(a \neq \{0\}\) , consider a non-zero element of a of minimal degree.)
\end{enumerate}

\begin{enumerate}
\def\labelenumi{\arabic{enumi}.}
\setcounter{enumi}{1}
\tightlist
\item
  Let g be a Lie algebra which is a free module; it is identified by means of \(\sigma: g \to Ug\) with its image in the enveloping algebra \(Ug\) . Let \(U^{+}g\) be the kernel of the canonical homomorphism \(Ug \to K\) . Let i be a mapping of X into g such that \(i(X)\) generates g as a Lie algebra.
\end{enumerate}

\begin{enumerate}
\def\labelenumi{(\alph{enumi})}
\item
  Show that \(\mathbf{U}^{+}\mathfrak{g}\) is generated by \(i(\mathbf{X})\) as a left Ug-module.
\item
  Show the equivalence of the following properties:
\item
  g is free with basic family i: X → g.
\end{enumerate}

\begin{enumerate}
\def\labelenumi{(\roman{enumi})}
\setcounter{enumi}{1}
\tightlist
\item
  The family \(i\) is a basis of the left Ug-module \(\mathbf{U}^{+}\mathfrak{g}\) .
\end{enumerate}

(The implication (i) \(\Rightarrow\) (ii) follows from Exercise 1 (a) using the isomorphism \(\mathrm{Ug} \to \mathrm{A}(\mathrm{X})\) . To show (ii) \(\Rightarrow\) (i), apply Exercise 1 (b) to the kernel of the homomorphism \(\mathrm{A}(\mathrm{X}) \to \mathrm{Ug}\) defined by \(i\) .)

\begin{enumerate}
\def\labelenumi{\arabic{enumi}.}
\setcounter{enumi}{2}
\tightlist
\item
  Let \(x \in \mathbf{X}\) . Show that the centralizer of \(x\) in \(\mathrm{A}(\mathbf{X})\) is the subalgebra generated by \(x\) . Deduce that the only elements of \(\mathrm{L}(\mathbf{X})\) which commute with \(x\) are the multiples of \(x\) . In particular, the centre of \(\mathrm{L}(\mathbf{X})\) is (0) if \(\operatorname{Card}(\mathbf{X}) \geqslant 2\) and the centre of \(\mathrm{L}(\mathbf{X}) / \left( \sum_{n > p} \mathrm{L}^n(\mathbf{X}) \right)\) reduces to the canonical image of \(\mathrm{L}^p(\mathbf{X})\) .
\end{enumerate}


¶ 4. Suppose that K is a field of characteristic p \textgreater{} 0.

\begin{enumerate}
\def\labelenumi{(\alph{enumi})}
\item
  Let \(\mathbf{H}\) be a basis of the Lie subalgebra \(\mathbf{L}(\mathbf{X})\) of \(\mathbf{A}(\mathbf{X})\) . Show (using the isomorphism \(\mathbf{U}(\mathbf{L}(\mathbf{X})) \to \mathbf{A}(\mathbf{X})\) and the Poincaré-Birkhoff-Witt Theorem) that the elements \(h^{p^n}\) ( \(h \in \mathbf{H}\) , \(n \in \mathbf{N}\) ) are linearly independent over \(\mathbf{K}\) .
\item
  If \(m\) is an integer \(\geqslant 0\) , let \(L_m\) denote the submodule of \(A(X)\) with basis the \(h^{p^n}\) , where \(h \in H, n \leqslant m\) . Show that \(L_m\) is a Lie subalgebra of \(A(X)\) and that, if \(a \in L_m - 1\) , then \(a^p \in L_m\) (argue by induction on \(m\) and use the Jacobson formulae, cf.~Chapter I, § 1, Exercise 19). Deduce that \(L_m\) does not depend on the choice of \(H\) and that it is the Lie subalgebra of \(A(X)\) generated by the \(x^{p^n}\) , where \(x \in X, n \leqslant m\) .
\item
  Let \(\mathbf{L}(\mathbf{X}, p)\) be the union of the \(\mathbf{L}_m\) for \(m \geqslant 0\) . Show that \(\mathbf{L}(\mathbf{X}, p)\) is the smallest Lie \(p\) -subalgebra of \(\mathbf{A}(\mathbf{X})\) containing \(\mathbf{X}\) (cf.~Chapter I, §1, Exercise 20); it admits as basis the family of \(h^{p^n}\) , where \(h \in \mathbf{H}\) , \(n \in \mathbf{N}\) .
\item
  Let \(\tilde{\mathbf{U}}(\mathbf{L}(\mathbf{X},p))\) be the restricted enveloping algebra of \(\mathbf{L}(\mathbf{X},p)\) , cf.~Chapter I, § 2, Exercise 6. Show that the injection of \(\mathbf{L}(\mathbf{X},p)\) into \(\mathbf{A}(\mathbf{X})\) can be extended to an isomorphism of \(\tilde{\mathbf{U}}(\mathbf{L}(\mathbf{X},p))\) onto \(\mathbf{A}(\mathbf{X})\) . (Use the Poincaré-Birkhoff-Witt Theorem and the exercise quoted above.) Deduce that \(\mathbf{L}(\mathbf{X},p)\) is the set of primitive elements of \(\mathbf{A}(\mathbf{X})\) , cf.~§ 1, Exercise 12.
\item
  Let \(f\) be a mapping of \(\mathbf{X}\) into a Lie \(p\) -algebra \(\mathfrak{g}\) . Show that \(f\) can be extended uniquely to a \(p\) -homomorphism \(\mathbf{F} \colon \mathbf{L}(\mathbf{X}, p) \to \mathfrak{g}\) . (Begin by extending \(f\) to an algebra homomorphism of \(\mathrm{A}(\mathbf{X})\) into \(\tilde{\mathrm{U}}\mathfrak{g}\) .)
\end{enumerate}

(The Lie \(p\) -algebra \(\mathbf{L}(\mathbf{X}, p)\) is called the free Lie \(p\) -algebra on the set \(\mathbf{X}\) .)

\exercisesfor{4}

In the following exercises the letter G denotes a group.

\begin{enumerate}
\def\labelenumi{\arabic{enumi}.}
\item
  Let \((\mathbf{G}_n)\) be an integral central filtration on \(G\) . Show that the Lie algebra \(\mathrm{gr}(G)\) is generated by \(\mathrm{gr}_1(G)\) if and only if \(\mathbf{G}_n = \mathbf{G}_{n+1}.\mathbf{C}^n\mathbf{G}\) for all \(n \geqslant 1\) . In that case, show that \(\mathbf{G}_n = \mathbf{G}_m.\mathbf{C}^n\mathbf{G}\) for \(m > n\) and deduce that \((\mathbf{G}_n) = (\mathbf{C}^n\mathbf{G})\) if there exists an integer \(m\) such that \(\mathbf{G}_m = \{\varepsilon\}\) .
\item
  Let \((\mathbf{G}_{\alpha})\) be a real filtration on a group \(\mathbf{G}\) and let \(v\) be the corresponding order function. Let \(\mathbf{H}\) be a subgroup of \(\mathbf{G}\) .
\end{enumerate}

\begin{enumerate}
\def\labelenumi{(\alph{enumi})}
\item
  Let \(H_{\alpha} = H \cap G_{\alpha}\) . Show that \((\mathrm{H}_{\alpha})\) is a real filtration of H and that the corresponding order function is the restriction of v to H. \(H_{\alpha}^{+} = H \cap G_{\alpha}^{+}\) and \(\mathrm{gr}(\mathrm{H})\) is identified with a graded subgroup of \(\mathrm{gr}(\mathrm{G})\) .
\item
  Suppose that H is normal and that \(v(G) \cap \mathbf{R}\) is a discrete subset of R. We write \((G/H)_{\alpha} = (G_{\alpha}H)/H\) . Show that \(((G/H)_{\alpha})\) is a real filtration of G/H and that the corresponding order function \(v_{G/H}\) is given by the formula
\end{enumerate}

\[
v _ {\mathrm{G/H}} (x) = \operatorname * {S u p} _ {y \in x} v (y).
\]

Show that, for all \(\alpha \in \mathbf{R}\) , there is an exact sequence

\[
0 \rightarrow \mathrm{gr} _ {\alpha} (\mathrm{H}) \rightarrow \mathrm{gr} _ {\alpha} (\mathrm{G}) \rightarrow \mathrm{gr} _ {\alpha} (\mathrm{G} / \mathrm{H}) \rightarrow 0.
\]

\begin{enumerate}
\def\labelenumi{(\alph{enumi})}
\setcounter{enumi}{2}
\tightlist
\item
  With the hypotheses of \((b)\) , suppose that \((G_{\alpha})\) is a central filtration. Show that so are the filtrations on H and G/H induced by \((G_{\alpha})\) and that the Lie algebra \(\mathrm{gr}(\mathrm{G}/\mathrm{H})\) is identified with the quotient of \(\mathrm{gr}(\mathrm{G})\) by the ideal \(\mathrm{gr}(\mathrm{H})\) .
\end{enumerate}

¶ 3. Let \((\mathrm{H}_{\alpha})\) be a real filtration on a group H and let \(v_{\mathrm{H}}\) be the corresponding order function. Suppose that G operates on H on the left and that, for all \(g \in G\) , the mapping \(h \mapsto g(h)\) is an automorphism of the filtered group H.

\begin{enumerate}
\def\labelenumi{(\alph{enumi})}
\tightlist
\item
  If \(\alpha \in \mathbf{R}\) , we denote by \(G_{\alpha}\) the set of \(g \in G\) such that
\end{enumerate}

\[
v _ {\mathrm{H}} (h ^ {- 1} g (h)) \geqslant v _ {\mathrm{H}} (h) + \alpha \quad \text { for   all } h \in \mathrm{H}.
\]

Show that \((\mathbf{G}_{\alpha})\) is a real filtration of \(\mathbf{G}\) and that \(\mathbf{G}_{\alpha} = \mathbf{G}\) if \(\alpha \leqslant 0\) . Let \(v\) denote the corresponding order function.

\begin{enumerate}
\def\labelenumi{(\alph{enumi})}
\setcounter{enumi}{1}
\tightlist
\item
  Suppose that the filtration \((\mathrm{H}_{\alpha})\) is central and that \(\mathrm{G}_0^+ = \mathrm{G}\) . Show that \((\mathrm{G}_{\alpha})\) is a central filtration of G. If \(\xi \in \mathrm{gr}_{\alpha}(\mathrm{G})\) and \(\eta \in \mathrm{gr}_{\beta}(\mathrm{H})\) , let \(g\) (resp. \(h\) ) be a representative of \(\xi\) (resp. \(\eta\) ) in \(\mathrm{G}_{\alpha}\) (resp. \(\mathrm{H}_{\beta}\) ); show that the image of \(h^{-1}g(h)\) in \(\mathrm{gr}_{\alpha + \beta}(\mathrm{H})\) is independent of the choice of \(g\) and \(h\) ; if it is denoted by \(\mathrm{D}_{\xi}(\eta)\) , show that \(\mathrm{D}_{\xi}\) can be extended to a derivation of the Lie algebra \(\mathrm{gr}(\mathrm{H})\) of degree \(\alpha\) and that \(\xi \mapsto \mathrm{D}_{\xi}\) is a homomorphism of the Lie algebra \(\mathrm{gr}(\mathrm{G})\) into the Lie algebra of derivations of \(\mathrm{gr}(\mathrm{H})\) . If \(v_{\mathrm{H}}(\mathrm{H}) \cap \mathbf{R}\) is contained in a discrete subgroup \(\Gamma\) of \(\mathbf{R}\) , so is \(v(\mathrm{G}) \cap \mathbf{R}\) and, for all \(\alpha \in \mathbf{R}\) , the mapping \(\xi \mapsto \mathrm{D}_{\xi}\) defined above is injective.
\end{enumerate}

\begin{enumerate}
\def\labelenumi{\arabic{enumi}.}
\setcounter{enumi}{3}
\item
  Let H be a nilpotent group of class c and let G be the group of automorphisms of H which act trivially on \(\mathrm{H}/(\mathrm{H},\mathrm{H})\) . Show that G is nilpotent of class \(\leqslant c - 1\) . (Apply Exercise 3 to the lower central series of H and note that G operates trivially on \(\mathrm{gr}(\mathrm{H})\) .) Show that, if H is a finite p-group (p prime), so is G (same method).
\item
  Let \(\mathbf{K}\) be a commutative field and let \(\mathbf{L}\) be a finite Galois extension of \(\mathbf{K}\) with Galois group \(\mathbf{G}\) . Let \(v_{\mathbf{L}}\) be a valuation of \(\mathbf{L}\) with values in \(\mathbf{Z}\) (Commutative Algebra, Chapter VI, § 3, no. 2) which is invariant under \(\mathbf{G}\) . If \(g \in \mathbf{G}\) , we write
\end{enumerate}

\[
v (g) = \sup _ {x \in L ^ {\bullet}} v _ {L} \left(\frac {g (x) - x}{x}\right).
\]

Show that v is the order function of a separated integral filtration \((G_{n})\) on G such that \(G_{0} = G\) and that the restriction of this filtration to \(G_{1}\) is central (apply Exercise 3, taking H to be the ideal of \(v_{L}\) ). Show that \(G_{1}\) is a p-group (resp. reduces to the identity element) if the residue field of L is of characteristic p \textgreater{} 0 (resp. of characteristic zero); when the fields L and K have the same residue field, \(G_{1}\) is the kernel of the homomorphism \(\phi\) defined in Commutative Algebra, Chapter VI, § 8, Exercise 11 (b).

¶ 6. Let K{[}G{]} be the algebra of G over K and let I be the kernel of the canonical homomorphism K{[}G{]} → K (``augmentation ideal''). Then


\(K[G] = K \oplus I\) and \(I\) admits as basis the family of \(g - 1\) where \(g \in G - \{e\}\) .

\begin{enumerate}
\def\labelenumi{(\alph{enumi})}
\item
  If \(n\) is an integer \(\geqslant 0\) , we denote by \(I^n\) the ideal of \(K[G]\) the \(n\) -th power of \(I\) . Let \(G_n\) be the set of \(g \in G\) such that \(g - 1 \in I^n\) . Show that \((G_n)\) is an integral central filtration on \(G\) . In particular, \(G_n \supset C^n G\) for all \(n\) .
\item
  Show that if \(\mathbf{K} = \mathbf{Z}\) the mapping \(g \mapsto g - 1\) defines on taking quotients an isomorphism of \(\mathrm{G} / (\mathrm{G}, \mathrm{G})\) onto \(\mathrm{I} / \mathrm{I}^2\) . Deduce that \(\mathrm{G}_2 = \mathrm{C}^2\mathrm{G}.\dagger\)
\item
  Suppose that K is a field of characteristic zero and that G is finite. Show that \(I^{n} = I\) for all \(n \geqslant 1\) .
\item
  Suppose that K is a field of characteristic p \textgreater{} 0 and that G is a p-group. Show that \(I^{n}\{0\}\) for n sufficiently large. (Show first, using Algebra, Chapter I, § 6, no. 5, Proposition 11, that every simple K{[}G{]}-module is isomorphic to K; deduce that I is the radical of K{[}G{]} and hence is nilpotent since K{[}G{]} is of finite rank over K.)
\item
  Suppose that K = Z and that G has the following property: for all \(g \in G\) such that \(g \neq 1\) , there exists a prime number p, a p-group P and a homomorphism \(f: G \to P\) such that \(f(g) \neq e\) . Show that in that case \(\bigcap_{n} I^{n} = \{0\}\) . (Reduce it immediately to the case where G = P. By applying (d) to the field \(F_{p}\) , we see that there exists m such that \(I^{m} \subset p\) . \(Z[G]\) and, as I is a direct factor in Z{[}G{]}, this implies that \(I^{m} \subset pI\) , whence \(\bigcap_{n} I^{mn} \subset \bigcap_{n} p^{n}I\) , which reduces to 0 since I is a finitely generated Abelian group.)
\end{enumerate}

\begin{enumerate}
\def\labelenumi{\arabic{enumi}.}
\setcounter{enumi}{6}
\item
  Let G be given the filtration \((\mathrm{C}^{\mathrm{n}}\mathrm{G})\) and suppose that \(\mathrm{gr}_{1}(\mathrm{G}) = \mathrm{G}/(\mathrm{G}, \mathrm{G})\) is cyclic. Show that \(\mathrm{gr}_{n}(\mathrm{G}) = \{0\}\) for \(n \geqslant 2\) (use Proposition 5) and deduce that \(\mathrm{C}^{\mathrm{n}}\mathrm{G} = (\mathrm{G}, \mathrm{G})\) for \(n \geqslant 2\) .
\item
  Let \(G = \mathbf{S}\mathbf{L}_{2}(\mathbf{Z})\) and let \(x = \begin{pmatrix} 1 & 1 \\ 0 & 1 \end{pmatrix}, y = \begin{pmatrix} 1 & 0 \\ 1 & 1 \end{pmatrix}, w = \begin{pmatrix} 0 & 1 \\ -1 & 0 \end{pmatrix}.\)
\end{enumerate}

\begin{enumerate}
\def\labelenumi{(\alph{enumi})}
\item
  Verify the formulae \(w^4 = 1\) , \(w = xy^{-1}x\) , \(wxw^{-1} = y^{-1}\) .
\item
  If \(g = \begin{pmatrix} a & b \\ c & d \end{pmatrix}\) is an element of \(G\) , we write \(l(g) = |a| + |b|\) . Show that \(l(g) = 1\) if and only if \(g\) is of the form \(y^n w^a\) , where \(n \in \mathbf{Z}\) and \(0 \leqslant a \leqslant 3\) . If \(l(g) \geqslant 2\) , show that there exists a power \(h\) of \(x\) or \(y\) such that \(l(gh) < l(g)\) . Deduce that \(G\) is generated by \(\{x, y\}\) .
\item
  Using (a) and (b), show that G/(G, G) is generated by the image \(\xi\) of x and that \(\xi^{12} = e.\ddagger\) Deduce that \(C^{n}G = (G, G)\) for \(n \geqslant 2\) . (apply Exercise 7).
\end{enumerate}

\begin{enumerate}
\def\labelenumi{\arabic{enumi}.}
\setcounter{enumi}{8}
\tightlist
\item
  Let G be given the filtration \((C^{n}G)\).
\end{enumerate}

\begin{enumerate}
\def\labelenumi{(\alph{enumi})}
\item
  Show that the ideal of \(\operatorname{gr}(\mathbf{G})\) generated by \(\operatorname{gr}_2(\mathbf{G})\) is \(\sum_{q\geqslant 2}\operatorname{gr}_q(\mathbf{G})\) .
\item
  Let \((x_{i})_{i\in I}\) be a generating family of \(G\) and let \(m\geqslant 1\) . Suppose that, for all \((i,j)\in \mathrm{I}^2\) , \((x_i,x_j)^m\in \mathrm{C}^3\mathrm{G}\) . Show that, for all \(q\geqslant 2\) and all \(u\in \mathrm{C}^q\mathrm{G}\) , \(u^{m}\in \mathrm{C}^{q + 1}\mathrm{G}\) (use \((a)\) ).
\end{enumerate}

\begin{enumerate}
\def\labelenumi{\arabic{enumi}.}
\setcounter{enumi}{9}
\tightlist
\item
  Let \(x, y\) be in \(G\) and let \(r, s\) be two integers \(\geqslant 1\) . Suppose that \((x^r, y^s) = e\) . (a) Show that, for all \(n \geqslant 1\) , \((x, y)^{(rs)n} \in C^{n+2}G\) . (It can be assumed that \(G\) is generated by \(\{x, y\}\) ; then apply Exercise 9 (b), noting that \((x^r, y^s) \equiv (x, y)^{rs} \mod C^3G\) .)
\end{enumerate}

\begin{enumerate}
\def\labelenumi{(\alph{enumi})}
\setcounter{enumi}{1}
\tightlist
\item
  Suppose that \(x^r = y^s = e\) ; let \(t\) denote the g.c.d. of \(r\) and \(s\) . Show that \((x, y)^t \in \mathbf{C}^3\mathbf{G}\) and deduce that \((x, y)^{t^n} \in \mathbf{C}^{n+2}\mathbf{G}\) for all \(n \geqslant 1\) (same method).
\end{enumerate}

\begin{enumerate}
\def\labelenumi{\arabic{enumi}.}
\setcounter{enumi}{10}
\tightlist
\item
  Let H be a subgroup of G and let m be an integer \(\geqslant 1\) . Suppose that G is generated by a family \((x_{i})_{i\in I}\) such that \(x_{i}^{m}\in H\) for all i.
\end{enumerate}

\begin{enumerate}
\def\labelenumi{(\alph{enumi})}
\tightlist
\item
  Let H be given the filtration induced by the filtration (C \(^{n}\) G) on G and let gr(H) be identified with a graded Lie subalgebra of gr(G), cf.~Exercise 2. Show that, for all \(n \geqslant 0\) ,
\end{enumerate}

\[
m ^ {n}. \mathrm{gr} _ {n} (\mathrm{G}) \subset \mathrm{gr} _ {n} (\mathrm{H}).
\]

Deduce that, for all \(z \in \mathbf{H}.\mathbf{C}^n\mathbf{G}\) , \(z^m \in \mathbf{H}.\mathbf{C}^{n+1}\mathbf{G}\) .

\begin{enumerate}
\def\labelenumi{(\alph{enumi})}
\setcounter{enumi}{1}
\tightlist
\item
  If G is nilpotent, show that there exists an integer \(N \geqslant 0\) (depending only on the nilpotency class of G) such that \(z^{mN} \in H\) for all \(z \in G\) .
\end{enumerate}

\begin{enumerate}
\def\labelenumi{\arabic{enumi}.}
\setcounter{enumi}{11}
\item
  Suppose that G is nilpotent. Let H be a subgroup of G and let m be an integer \(\geqslant 1\) and let x, y be elements of G such that \(x^{m} \in H\) and \(y^{m} \in H\) . Show that there exists an integer \(N \geqslant 0\) such that \((xy)^{nN} \in H\) . (Apply Exercise 11 to the group generated by \(\{x, y\}\) and its intersection with H.)
\item
  \begin{enumerate}
  \def\labelenumii{(\alph{enumii})}
  \tightlist
  \item
    Let \(\mathbf{F}\) be a free group with basic family \(\{\overline{x},\overline{y}\}\) of two elements, let \(c\) be an integer \(\geqslant 2\) and let \(m\) be an integer \(\geqslant 1\) . We write \(\mathrm{F}^c = \mathrm{F} / C^c\mathrm{F}\) and denote by \(x,y\) the images of \(\overline{x},\overline{y}\) in \(\mathrm{F}^c\) . Let \(\mathrm{F}_m^c\) be the subgroup of \(\mathrm{F}^c\) generated by \(\{x^{m},y\}\) . Show that there exists an integer \(\mathrm{N}\geqslant 0\) such that \(z^{mN}\in \mathrm{F}_{m}^{c}\) for all \(z\in \mathrm{F}^c\) (use Exercise 11).
  \end{enumerate}
\end{enumerate}

\begin{enumerate}
\def\labelenumi{(\alph{enumi})}
\setcounter{enumi}{1}
\item
  Let \(\mathrm{I}^c\) (resp. \(\mathbf{I}_m^c\) ) be the normal subgroup of \(\mathrm{F}^c\) (resp. \(\mathbf{F}_m^c\) ) generated by \(y\) . Show that, if \(N\) is chosen as above and \(z \in \mathrm{I}^c\) , then \(z^{n^k} \in \mathbf{I}_m^c\) (note that \(\mathrm{F}^c / \mathrm{I}^c\) is an infinite cyclic group with generator the image of \(x\) and deduce that \(\mathbf{F}_m^c / \mathbf{I}_m^c \to \mathbf{F}^c / \mathrm{I}^c\) is injective and hence that \(\mathbf{I}_m^c = \mathbf{F}_m^c \cap \mathbf{I}^c\) ). In particular \(xy^m x^{-1} \in \mathbf{I}_m^c\) .
\item
  Suppose that G is nilpotent. Let H be a subgroup of G and let L be a normal subgroup of H. Let \(g \in G\) and let \(m \geqslant 1\) be such that \(g^{m} \in H\) . Show that, if N is sufficiently large, \(gl^{m^{N}} g^{-1} \in L\) for all \(l \in L\) . (If G is of class \textless c, choose N as in (a) and use the homomorphism \(f: F^{c} \to G\) such that \(f(x) = g, f(y) = l\) ; note that \(f(F_{m}^{c}) \subset H, f(I_{m}^{c}) \subset L\) and apply (b) above.)
\end{enumerate}

¶ 14. Let P be a set of prime numbers. An integer n is called a P-integer if it is ≠0 and all its prime factors belong to P. An element \(x \in G\) is called a P-torsion element if there exists a P-integer n such that \(x^{n} = \epsilon\) ; G is called a P-torsion (resp. P-torsion free) group if every element of G is a P-torsion element (resp. if no element of G other than e is a P-torsion element). G is called P-divisible if, for all \(x \in G\) and every P-integer n, there exists \(y \in G\) such that \(x = y^{n}\) .

Suppose that G is nilpotent.

\begin{enumerate}
\def\labelenumi{(\alph{enumi})}
\item
  Let H be a subgroup of G and let \(H_{P}\) be the set of \(x \in G\) such that there exists a P-integer n for which \(x^{n} \in H\) . Show that \(H_{P}\) is a subgroup of G (use Exercise 12) and that \((\mathrm{H}_{\mathrm{P}})_{\mathrm{P}} = \mathrm{H}_{\mathrm{P}}\) . The group \(H_{P}\) is called the P-saturation of H in G. If G is P-divisible, so is \(H_{P}\) .
\item
  Let \(\mathbf{L}\) be a normal subgroup of \(\mathbf{H}\) . Show that \(\mathbf{L}_{\mathbb{P}}\) is normal in \(\mathbf{H}_{\mathbb{P}}\) . (Use Exercise 13 (c) to prove that, if \(g \in \mathbf{H}_{\mathbb{P}}\) and \(l \in \mathbf{L}\) , then \(glg^{-1} \in \mathbf{L}_{\mathbb{P}}\) .)
\end{enumerate}

In particular, if L is normal in G, so is \(L_{P}\) and \(G/L_{P}\) is P-torsion-free. Deduce that the set of P-torsion elements of G is the smallest normal subgroup N of G such that G/N is P-torsion-free.

\begin{enumerate}
\def\labelenumi{(\alph{enumi})}
\setcounter{enumi}{2}
\item
  Suppose that G is P-torsion-free. Let \(n\) be a P-integer and let \(x, y \in G\) be such that \(x^n = y^n\) . Show that \(x = y\) . (Apply Exercise 10 (a) with \(r = s = n\) . Deduce that there exists a P-integer N such that \((x, y)^{\mathbb{N}} = e\) , whence \((x, y) = e\) since G is P-torsion-free; then \((x^{-1}y)^n = e\) , whence finally \(x = y\) .)
\item
  Let H be a subgroup of G, let L be a nilpotent group and let \(f: H \to L\) be a homomorphism. Let \(\Gamma\) be the graph of f in \(H \times L\) and let \(\Gamma_{P}\) be its P-saturation in \(G \times L\) . Show that \(\Gamma_{P}\) is contained in \(H_{P} \times L\) and that \(pr_{1}: \Gamma_{P} \to H_{P}\) is surjective if L is P-divisible and injective if L is P-torsion-free.
\end{enumerate}

Suppose that L is P-divisible and P-torsion-free. Show that f can be extended uniquely to a homomorphism \(f_{P}: H_{P} \to L\) and that the graph of \(f_{P}\) is \(\Gamma_{P}\) .

¶ 15. Preserving the notation of the above exercise, suppose that G is nilpotent. Let \(i: \mathbf{G} \to \overline{\mathbf{G}}\) be a homomorphism of G into a nilpotent group \(\overline{\mathbf{G}}\) . \((i, \overline{\mathbf{G}})\) is called a P-envelope of G if the following conditions hold:

\begin{enumerate}
\def\labelenumi{(\roman{enumi})}
\item
  \(\overline{\mathbf{G}}\) is P-divisible and P-torsion-free.
\item
  The kernel of \(i\) is the set of P-torsion elements of G.
\item
  The P-saturation of \(i(\mathbf{G})\) in \(\overline{\mathbf{G}}\) is equal to \(\overline{\mathbf{G}}\) .
\end{enumerate}

\begin{enumerate}
\def\labelenumi{(\alph{enumi})}
\item
  Let \((i, \overline{\mathbf{G}})\) be a P-envelope of G and let L be a P-divisible P-torsion-free nilpotent group. Show that, for every homomorphism \(f\colon \mathbf{G} \to \mathbf{L}\) , there exists one and only one homomorphism \(\bar{f}\colon \overline{\mathbf{G}} \to \mathbf{L}\) such that \(\bar{f} \circ i = f\) (reduce it to the case where G is P-torsion-free and use Exercise 14 (d)). Deduce that, if G has a P-envelope, \(\dagger\) this is unique up to isomorphism.
\item
  Let \((i, \overline{\mathbf{G}})\) be a P-envelope of G, let H be a subgroup of G, let \(\overline{\mathbf{H}}\) be the P-envelope of \(i(\mathbf{H})\) in \(\overline{\mathbf{G}}\) and let \(i_{\mathbf{H}}: \mathbf{H} \to \overline{\mathbf{H}}\) the homomorphism induced by \(i\) . Show that \((i_{\mathbf{H}}, \overline{\mathbf{H}})\) is a P-envelope of H.
\end{enumerate}

Suppose that H is normal in G. Then \(\bar{H}\) is normal in \(\bar{G}\) , cf.~Exercise 14 (b); if \(i_{G/H}: G/H \to \overline{G}/\overline{H}\) is the homomorphism induced by i, show that \((i_{G/H}, \overline{G}/\overline{H})\) is a P-envelope of G/H.

\begin{enumerate}
\def\labelenumi{\arabic{enumi}.}
\setcounter{enumi}{15}
\tightlist
\item
  We preserve the notation of the two previous exercises.
\end{enumerate}

\begin{enumerate}
\def\labelenumi{(\alph{enumi})}
\item
  Let \(\mathbf{S}_{\mathbf{P}}\) be the set of P-integers and let \(\mathbf{\dot{Z}}_{\mathbf{P}} = \mathbf{S}_{\mathbf{P}}^{-1}\mathbf{Z}\) be the ring of fractions of \(\mathbf{Z}\) defined by \(\mathbf{S}_{\mathbf{P}}\) (Commutative Algebra, Chapter II, § 2, no. 1). Suppose that G is nilpotent, P-torsion-free and P-divisible and let \(t \in \mathbf{Z}_{\mathbf{P}}\) , \(g \in G\) . Show that there exists one and only one element \(h\) of G such that \(h^s = g^{st}\) for all \(s \in \mathbf{S}_{\mathbf{P}}\) such that \(st \in \mathbf{Z}\) . The element \(h\) is called the \(t\) -th power of \(g\) and denoted by \(g^t\) . The mapping \(t \mapsto g^t\) is a homomorphism of \(\mathbf{Z}_{\mathbf{P}}\) into G. If \(t\) is invertible in \(\mathbf{Z}_{\mathbf{P}}\) , \(g \mapsto g^t\) is bijective. (Use Exercise 14 (c)).
\item
  Let \(\mathbf{A}\) be a commutative group and let \(i\) be the canonical mapping of \(\mathbf{A}\) into \(\mathbf{A}_{\mathbb{P}} = \mathbf{Z}_{\mathbb{P}} \times \mathbf{A}\) . Show that \((i, \mathbf{A}_{\mathbb{P}})\) is a P-envelope of \(\mathbf{A}\) .
\item
  Let \(\mathbf{G}\) (resp. \(\overline{\mathbf{G}}\) ) be the lower strict triangular group of order \(n\) over a ring \(k\) (resp. over the ring \(k_{\mathrm{P}} = k\otimes \mathbf{Z}_{\mathrm{P}}\) ) and let \(i\) be the homomorphism of \(\mathbf{G}\) into \(\overline{\mathbf{G}}\) defined by the canonical homomorphism of \(k\) into \(k_{\mathrm{P}}\) (Algebra, Chapter II, § 10, no. 4). Show that \((i,\overline{\mathbf{G}})\) is a P-envelope of \(\mathbf{G}\) .
\item
  Let G be a finite nilpotent group and hence the direct product of its Sylow \(p\) -groups \(G_p\) (Algebra, Chapter I, § 6, no. 7, Theorem 4). Let \(\overline{G}\) be the product of the \(G_p\) for \(p \notin P\) and let \(i\) be the canonical projection of G onto \(\overline{G}\) . Show that \((i, \overline{G})\) is a P-envelope of G.
\end{enumerate}

¶ 17. Suppose that G is nilpotent. Let P be a set of prime numbers and let \(\mathfrak{V}_{P}(G)\) be the set of normal subgroups of G of finite index whose index is a P-integer (cf.~Exercise 14). Let \(\mathcal{T}_{P}(G)\) be the topology defined on G by the filter base \(\mathfrak{V}_{P}(G)\) (General Topology, Chapter III, § 1, no. 2, Example).

\begin{enumerate}
\def\labelenumi{(\alph{enumi})}
\item
  Let N be a subgroup of G of finite index such that (G: N) is a P-integer and let \(N'\) be the intersection of the conjugates of N. Show that \((G: N')\) is a P-integer (use the fact that the group \(G/N'\) is the product of its Sylow groups); deduce that N is open with respect to \(\mathcal{T}_{\mathbb{P}}(G)\) .
\item
  Let H be a subgroup of G of finite index. Show that \(\mathcal{T}_{\mathrm{P}}(\mathrm{H})\) coincides with the topology induced on H by \(\mathcal{T}_{\mathrm{P}}(\mathrm{G})\) . (Same method.)
\item
  Let H be a normal subgroup of G such that G/H is isomorphic to Z and let x be a representative in G of a generator of G/H. Let \(N \in \mathfrak{V}_{\mathrm{P}}(\mathrm{H})\) and let \(N' = \bigcap_{n \in Z} x^{n} N x^{-1}\) . Show that \(N' \in \mathfrak{V}_{\mathrm{P}}(\mathrm{H})\) and that \(x N' x^{-1} = N'\) . If \(N''\) is the subgroup generated by \(N'\) and x, show that \((G: N'') = (H: N')\) . Deduce that \(\mathcal{T}_{\mathrm{P}}(\mathrm{H})\) is induced by \(\mathcal{T}_{\mathrm{P}}(\mathrm{G})\) .
\item
  Suppose that \(\mathbf{G}\) is finite. Show that, if \(\mathbf{H}\) is a subgroup of \(\mathbf{G}\) , there exists a sequence of subgroups
\end{enumerate}

\[
\mathrm{H} = \mathrm{H} _ {0} \subset \mathrm{H} _ {1} \subset \dots \subset \mathrm{H} _ {k} = \mathrm{G}
\]

such that \(\mathbf{H}_i\) is normal in \(\mathrm{H}_{i + 1}\) for \(0\leqslant i < k\) and that \(\mathrm{H}_{i + 1} / \mathrm{H}_i\) is finite or

isomorphic to \(\mathbf{Z}\) . Deduce, using \((b)\) and \((c)\) , that \(\mathcal{T}_{\mathrm{P}}(\mathrm{H})\) is induced by \(\mathcal{T}_{\mathrm{P}}(\mathrm{G})\) .

\begin{enumerate}
\def\labelenumi{(\alph{enumi})}
\setcounter{enumi}{4}
\tightlist
\item
  With the hypotheses of (d), let P' denote the complement of P in the set of prime numbers. Show that the closure of H under \(\mathcal{T}_{\mathrm{P}}(\mathrm{G})\) coincides with the P'-saturation of H in G (Exercise 14); in particular, G is Hausdorff if and only if it is P'-torsion-free.
\end{enumerate}

\begin{enumerate}
\def\labelenumi{\arabic{enumi}.}
\setcounter{enumi}{17}
\tightlist
\item
  The upper central series \((Z_{i}G)\) of the group G is defined inductively as follows:
\end{enumerate}

\begin{enumerate}
\def\labelenumi{(\roman{enumi})}
\item
  \(Z_{i}G = \{e\}\) if \(i \leqslant 0\) ;
\item
  \(Z_{i}G/Z_{i-1}G\) is the centre of \(G/Z_{i-1}G\) .
\end{enumerate}

Then \(\{e\} = Z_0G \subset Z_1G \subset \cdots\) and \(Z_1G\) is the centre of G; the \(Z_iG\) are characteristic subgroups of G.

\begin{enumerate}
\def\labelenumi{(\alph{enumi})}
\item
  Show that \(\mathbf{G}\) is nilpotent of class \(\leqslant c\) if and only if \(\mathbf{G} = \mathbf{Z}_c\mathbf{G}\) .
\item
  Show that \((\mathbf{C}^n\mathbf{G},\bar{\mathbf{Z}}_m\mathbf{G})\subset \mathbf{Z}_{m - n}\mathbf{G}.\)
\item
  Suppose that G is nilpotent and P-torsion-free (where P is a set of prime numbers, cf.~Exercise 14). Show that the \(Z_{1}G\) are P-saturated. (It suffices to verify this for \(Z_{1}G\) ; if n is a P-integer and \(g \in G\) is such that \(g^{n} \in Z_{1}G\) , then \(xg^{n}x^{-1} = g^{n}\) for all \(x \in G\) , whence \(xgx^{-1} = g\) by Exercise 14 (c) and it follows that g belongs to \(Z_{1}G\) .)
\end{enumerate}

\exercisesfor{5}

In the following exercises we assume the hypotheses and notation of §5. F denotes the free group \(\mathrm{F}(\mathbf{X})\) and g the unique homomorphism of F into the Magnus group \(\Gamma(\mathbf{X})\) such that \(g(x)=1+x\) for all \(x\in X\) (cf.~Theorem 1).

\begin{enumerate}
\def\labelenumi{\arabic{enumi}.}
\tightlist
\item
  Let the algebra \(\mathbf{K}[\mathbf{F}]\) of \(\mathbf{F}\) be given the filtration \((\mathbf{I}^n)\) consisting of the powers of the augmentation ideal \(\mathbf{I}\) ( \(\S 4\) , Exercise 6).
\end{enumerate}

\begin{enumerate}
\def\labelenumi{(\alph{enumi})}
\item
  Let \(\tilde{g}\colon \mathrm{K}[\mathrm{F}] \to \tilde{\mathrm{A}}(\mathrm{X})\) be the unique algebra homomorphism extending \(g\colon \mathrm{F} \to \tilde{\mathrm{A}}(\mathrm{X})^*\) . Show that \(\tilde{g}\) maps \(\mathrm{I}^{\mathfrak{n}}\) to the ideal \(\tilde{\mathrm{A}}_{\mathfrak{n}}(\mathrm{X})\) (cf.~no. 1) and defines when taking quotients an isomorphic \(\tilde{g}_{n}\) of \(\mathrm{K}[\mathrm{F}]/\mathrm{I}^{\mathfrak{n}}\) onto \(\tilde{\mathrm{A}}(\mathrm{X})/\tilde{\mathrm{A}}_{\mathfrak{n}}(\mathrm{X})\) . (Define an inverse homomorphism of \(\tilde{g}_{n}\) by means of the homomorphism of \(\mathrm{A}(\mathrm{X})\) into \(\mathrm{K}[\mathrm{F}]\) which maps \(x\) to \(x - 1\) for all \(x \in \mathrm{X}\) .) Deduce that \(\tilde{\mathrm{A}}(\mathrm{X})\) is isomorphic to the Hausdorff completion of \(\mathrm{K}[\mathrm{F}]\) under the topology defined by \((\mathrm{I}^{\mathfrak{n}})\) .
\item
  Suppose that \(\mathbf{K} = \mathbf{Z}\) . Show that the filtration \((\mathrm{I}^{\mathrm{n}})\) is separated (use Proposition 3 and Exercise 6 (e) of §4) and that the filtration of \(\mathbf{F}\) defined by it (Exercise 6 (a) of §4) coincides with the filtration \((\mathbf{C}^{\mathrm{n}}\mathbf{F})\) . Deduce that \(\tilde{g}\colon \mathbf{Z}[\mathrm{F}] \to \hat{\mathrm{A}}_{\mathbf{Z}}(\mathrm{X})\) is injective.
\end{enumerate}

\begin{enumerate}
\def\labelenumi{\arabic{enumi}.}
\setcounter{enumi}{1}
\tightlist
\item
  Let G be a group, let K{[}G{]} be its algebra over K and let I be its augmentation ideal (§ 4, Exercise 6). Let \(\varepsilon\) denote the canonical homomorphism of K{[}G{]} onto K; then \(\operatorname{Ker}(\varepsilon) = I\) .
\end{enumerate}

\begin{enumerate}
\def\labelenumi{(\alph{enumi})}
\tightlist
\item
  Let M be a left K{[}G{]}-module and let Z(G, M) be the group of crossed homomorphisms of G into M (Algebra, Chapter I, § 6, Exercise 7). If \(f \in \operatorname{Hom}_{\mathbb{K}[\mathbb{G}]}(\mathrm{I}, \mathrm{M})\) , let \(f_1\) be the mapping \(g \mapsto f(g - 1)\) of \(\mathbf{G}\) into \(\mathbf{M}\) . Show that \(f_1\) is a crossed homomorphism (use the identity
\end{enumerate}

\[
g g ^ {\prime} - 1 = g (g ^ {\prime} - 1) + g - 1
\]

to show that \(f_{1}(gg') = g \cdot f_{1}(g') + f_{1}(g)\) and that \(f \mapsto f_{1}\) is an isomorphism of \(\operatorname{Hom}_{\mathbb{K}[G]}(\mathbf{I}, \mathbf{M})\) onto \(\mathbf{Z}(\mathbf{G}, \mathbf{M})\) .

\begin{enumerate}
\def\labelenumi{(\alph{enumi})}
\setcounter{enumi}{1}
\item
  Take G to be the free group \(\mathrm{F} = \mathrm{F}(\mathrm{X})\) . Show that, for every mapping \(\theta: \mathrm{X} \to \mathrm{M}\) , there exists one and only one element \(f_{\theta}\) of \(\mathrm{Z}(\mathrm{M}, \mathrm{F})\) which extends \(\theta\) . (Use the interpretation of \(\mathrm{Z}(\mathrm{F}, \mathrm{M})\) in terms of the semi-direct product of F by M, cf.~Algebra, Chapter I, loc. cit.) Deduce, using (a), that the family \((x - 1)_{x \in \mathbb{X}}\) is a basis of I as a left K{[}F{]}-module.
\item
  By (b) every element \(u \in \mathbf{K}[\mathbf{F}]\) can be written uniquely in the form
\end{enumerate}

\[
u = \varepsilon (u) + \sum_ {x \in \mathbf {X}} \mathrm{D} _ {x} (u) (x - 1),
\]

where \(\mathrm{D}_{x}(u)\in\mathrm{K}[\mathrm{F}]\) and the \(\mathrm{D}_{x}(u)\) are zero except for a finite number. The mapping \(u\mapsto\mathrm{D}_{x}(u)\) is called the partial derivative with respect to x. It is characterized by the following properties:

\begin{enumerate}
\def\labelenumi{(\roman{enumi})}
\item
  \(D_{x}\) is a K-linear mapping of K{[}F{]} into K{[}F{]}.
\item
  \(\mathrm{D}_{x}(uv)=u.\mathrm{D}_{x}(v)+\varepsilon(v).\mathrm{D}_{x}(u)\) for u, v in K{[}F{]}.
\item
  \(\mathrm{D}_{x}(x)=1\) and \(\mathrm{D}_{x}(y)=0\) if \(y\in X-\{x\}\) .
\end{enumerate}

If \(n\geqslant 1\) , then

\[
\begin{array}{c} \mathrm{D} _ {x} (x ^ {n}) = 1 + x + \dots + x ^ {n - 1} \\ \mathrm{D} _ {x} (x ^ {- n}) = - x ^ {- n} - x ^ {- n + 1} - \dots - x ^ {- 1}. \end{array}
\]

\begin{enumerate}
\def\labelenumi{(\alph{enumi})}
\setcounter{enumi}{3}
\tightlist
\item
  Let \(\varepsilon_{\mathrm{A}}\) be the canonical homomorphism of A(X) onto K. Show that every element \(u\) of A(X) can be written uniquely in the form
\end{enumerate}

\[
u = \varepsilon_ {\mathrm{A}} (u) + \sum_ {x \in \mathbf {X}} d _ {x} (u). x,
\]

where \(d_{x}(u) \in \mathrm{A}(\mathrm{X})\) and the \(d_{x}(u)\) are zero except for a finite number. Show that \(d_{x}: u \to d_{x}(u)\) is an endomorphism of \(\mathrm{A}(\mathrm{X})\) which can be extended by continuity to \(\hat{\mathrm{A}}(\mathrm{X})\) and enjoys properties analogous to (i), (ii), (iii) above. Show \(d_{x} \circ \tilde{g} = \tilde{g} \circ D_{x}\) , where \(\tilde{g}\) is the homomorphism of K{[}F{]} into \(\hat{\mathrm{A}}(\mathrm{X})\) which extends g (cf.~Exercise 1).†

¶ 3. We preserve the notation of Exercise 2.

\begin{enumerate}
\def\labelenumi{(\alph{enumi})}
\tightlist
\item
  Let \(J\) be a right ideal of \(K[F]\) contained in \(I\) and let \(u\) be an element of \(I\) . Show that
\end{enumerate}

\[
u \in J. I \Leftrightarrow D _ {x} (u) \in J \quad \text { for   all } x \in X.
\]

In particular, an element of I belongs to \(\mathbf{I}^n\) ( \(n \geqslant 1\) ) if and only if all its partial derivatives belong to \(\mathbf{I}^{n-1}\) .


\begin{enumerate}
\def\labelenumi{(\alph{enumi})}
\setcounter{enumi}{1}
\tightlist
\item
  Let R be a normal subgroup of F, let G = F/R and let J be the kernel of the canonical homomorphism \(\gamma: K[F] \to K[G]\) . Show that J is generated as a left (resp. right) ideal by the elements r - 1, where \(r \in R\) . Prove the exactness of the sequence
\end{enumerate}

\[
0 \longrightarrow \mathrm{J} / (\mathrm{J}. \mathrm{I}) \stackrel {\alpha} {\longrightarrow} \mathrm{I} / (\mathrm{J}. \mathrm{I}) \stackrel {\beta} {\longrightarrow} \mathrm{K} [ \mathrm{G} ] \stackrel {\delta} {\longrightarrow} \mathrm{K} \longrightarrow 0,\tag{*}
\]

where \(\alpha\) is derived by taking quotients from the inclusion of J in I and \(\beta\) is derived by taking quotients from the restriction of \(\gamma\) to I.

\begin{enumerate}
\def\labelenumi{(\alph{enumi})}
\setcounter{enumi}{2}
\item
  For \(u \in \mathbf{I}\) , \(x \in \mathbf{X}\) , let \(\overline{\mathrm{D}}_x(u)\) denote the image of \(\mathrm{D}_x(u)\) in \(\mathrm{K}[\mathrm{G}]\) under \(\gamma\) . Show, using (a), that the family \((\overline{\mathrm{D}}_x)_{x \in \mathbf{X}}\) defines on taking quotients an isomorphism of \(\mathrm{I} / (\mathrm{J}. \mathrm{I})\) onto \(\mathrm{K}[\mathrm{G}]^{(\mathbb{X})}\) .
\item
  For \(r \in \mathbb{R}\) , let \(\theta(r)\) be the image of \(r - 1\) in \(J/(J.I)\) . Show that
\end{enumerate}

\[
\theta (r r ^ {\prime}) = \theta (r) + \theta (r ^ {\prime}) \quad \text { and } \quad \theta (y r y ^ {- 1}) = y. \theta (r) \quad \text { for   } r, r ^ {\prime} \text {   in   } R, y \in F.
\]

(Use the identities

\[
\begin{array}{c} {r r ^ {\prime} - 1 = (r - 1) (r ^ {\prime} - 1) + (r - 1) + (r ^ {\prime} - 1)} \\ {y r y ^ {- 1} - 1 = y (r - 1) (y ^ {- 1} - 1) + y (r - 1).)} \end{array}
\]

Deduce that \(\theta\) defines a homomorphism of \(R/(R, R)\) into \(J/(J, I)\) compatible with the action of \(G\) and show that the image of this homomorphism generates the K-module \(J/(J, I)\) ; when \(K = Z\) , show that an isomorphism is thus obtained of \(R/(R, R)\) onto \(J/(J, I)\) (define the inverse homomorphism directly). (e) Let \((r_{\alpha})_{\alpha \in A}\) be a family of elements of \(R\) generating \(R\) as a normal subgroup of \(F\) . Show that the \(\theta(r_{\alpha})\) generate the K{[}G{]}-module \(J/(J, I)\) .

\begin{enumerate}
\def\labelenumi{(\alph{enumi})}
\setcounter{enumi}{5}
\tightlist
\item
  The matrix \((\overline{\mathrm{D}}_x(r_\alpha))_{x\in \mathbf{X},\alpha \in \mathbf{A}}\) defines a homomorphism
\end{enumerate}

\[
\rho \colon \mathrm{K} [ \mathrm{G} ] ^ {(\mathrm{A})} \rightarrow \mathrm{K} [ \mathrm{G} ] ^ {(\mathrm{X})}
\]

of left K{[}G{]}-modules. Show that the sequence

\[
\mathrm{K} [ \mathrm{G} ] ^ {(\mathrm{A})} \xrightarrow {\rho} \mathrm{K} [ \mathrm{G} ] ^ {(\mathrm{X})} \xrightarrow {\delta} \mathrm{K} [ \mathrm{G} ] \xrightarrow {\varepsilon} \mathrm{K} \longrightarrow 0\tag{**}
\]

is exact, where \(\delta\) is the homomorphism \((u_x) \mapsto \sum_{x \in X} u_x(\gamma(x) - 1)\) .

(Transform the exact sequence (*) using (c), (d), (e) above.)†

\begin{enumerate}
\def\labelenumi{\arabic{enumi}.}
\setcounter{enumi}{3}
\tightlist
\item
  For all \(n \in \mathbf{N}\) , let \(\binom{\mathrm{T}}{n}\) denote the polynomial
\end{enumerate}

\[
\mathrm{T} (\mathrm{T} - 1) \dots (\mathrm{T} - n + 1) / n!.
\]

If \(f(\mathrm{T})\) is a polynomial, let \(\Delta f\) denote the polynomial \(f(\mathrm{T} + 1) - f(\mathrm{T})\) .

Then \(\Delta\binom{\mathrm{T}}{0} = \Delta 1 = 0\) and \(\Delta\binom{\mathrm{T}}{n} = \binom{\mathrm{T}}{n - 1}\) if \(n \geqslant 1\) .

\begin{enumerate}
\def\labelenumi{(\alph{enumi})}
\item
  Let \(f \in \mathbf{Q}[\mathrm{T}]\) . Show the equivalence of the following properties:
\item
  \(f\) maps \(\mathbf{Z}\) into itself.
\end{enumerate}

\begin{enumerate}
\def\labelenumi{(\roman{enumi})}
\setcounter{enumi}{1}
\item
  \(f\) is a linear combination with coefficients in \(\mathbf{Z}\) of the \(\binom{\mathrm{T}}{n}\) .
\item
  \(f(0)\in \mathbf{Z}\) and \(\Delta f\) maps \(\mathbf{Z}\) into itself.
\end{enumerate}

(Argue by induction on \(\deg(f)\) , noting that \(\deg(\Delta f) = \deg(f) - 1\) if \(\deg(f) \neq 0\) .)

A polynomial \(f\) satisfying the above properties is called binomial. The sum, product and composition of two binomial polynomials is a binomial polynomial.

\begin{enumerate}
\def\labelenumi{(\alph{enumi})}
\setcounter{enumi}{1}
\item
  Let \(p\) be a prime number and let \(f\) be a binomial polynomial. Show that \(f\) maps the ring \(\mathbf{Z}_p\) of \(p\) -adic integers into itself (use the continuity of \(f\) and the fact that \(\mathbf{Z}\) is dense in \(\mathbf{Z}_p\) ).
\item
  Let \(\mathbf{P}\) be a set of prime numbers, let \(\mathrm{S}_P\) be the set of P-integers (§ 4, Exercise 14) and let \(\mathbf{Z}_P = \mathrm{S}_{\mathrm{P}}^{-1}\mathbf{Z}\) . Show that if \(f\) is a binomial polynomial, \(f\) maps \(\mathbf{Z}_{\mathrm{P}}\) into itself (apply (b) to the prime numbers \(p\) not belonging to \(\mathbf{P}\) ).
\end{enumerate}

\begin{enumerate}
\def\labelenumi{\arabic{enumi}.}
\setcounter{enumi}{4}
\tightlist
\item
  Let \(\Gamma = \Gamma(X)\) be the Magnus group on \(X\) and let \((\Gamma_n)\) be its natural filtration (no. 2). Show that the graded Lie algebra \(\mathrm{gr}(\Gamma)\) corresponding to this filtration is isomorphic to the Lie subalgebra \(\bigoplus_{n\geqslant 1}A_n(X)\) of \(A(X)\) . If \(X\neq \emptyset\) , this Lie algebra is not generated by its elements of degree 1; deduce that \((\Gamma_n)\) is not the lower central series of \(\Gamma\) .
\end{enumerate}

¶ 6. Let P be a set of prime numbers, let \(S_P\) be the set of P-integers (§ 4, Exercise 14) and let \(Z_P = S_P^{-1}Z\) .

\begin{enumerate}
\def\labelenumi{(\alph{enumi})}
\tightlist
\item
  Suppose that, for all \(s \in S_{P}\) , the mapping \(k \mapsto sk\) is a bijection of K onto itself; this is equivalent to saying that K can be given a \(Z_{P}\) -algebra structure.
\end{enumerate}

With the notation of Exercise 5, show that, for all \(s \in S_{\mathbb{P}}\) , the mapping \(a \mapsto a^s\) of \(\Gamma\) into itself is bijective and that the same is true for each quotient \(\Gamma / \Gamma_n\) for \(n \geqslant 1\) . If \(t \in \mathbf{Z}_{\mathbb{P}}\) and \(a \in \Gamma\) (resp. \(a \in \Gamma / \Gamma_n\) ), \(a^s\) is defined as in Exercise 16 of §4; writing \(a\) in the form \(1 + \alpha\) , where \(\alpha \in \hat{\mathrm{A}}_1(\mathrm{X})\) , show that

\[
a ^ {t} = (1 + \alpha) ^ {t} = \sum_ {n = 0} ^ {\infty} \binom {t} {n} \alpha^ {n}.
\]

(Note that the coefficients \(\binom{t}{n}\) belong to \(\mathbf{Z}_{\mathrm{P}}\) , cf.~Exercise 4.)

\begin{enumerate}
\def\labelenumi{(\alph{enumi})}
\setcounter{enumi}{1}
\item
  Suppose now that \(\mathbf{K} = \mathbf{Z}_{\mathrm{P}}\) . Let \(c\) be an integer \(\geqslant 1\) . The group \(\mathbf{F}^c = \mathbf{F} / \mathbf{C}^c\mathbf{F}\) is identified with a subgroup of \(\Gamma /\Gamma_c\) by means of the homomorphism derived from \(g\) by taking quotients (notation of Theorem 2). The group \(\Gamma / \Gamma_c\) is a P-torsion-free P-divisible group of class \(< c\) . Let \(F_P^c\) be the P-saturation of \(F^c\) in \(\Gamma / \Gamma_c\) (§ 4, Exercise 14) and let \(i\) be the injection of \(F^c\) into \(F_P^c\) . The ordered pair \((i, F_P^c)\) is a P-envelope of \(F^c\) (§ 4, Exercise 15). Deduce that every nilpotent group has a P-envelope (note that every nilpotent group of class \(< c\) is the quotient of a group \(F^c\) for suitable X and use part (b) of Exercise 15 of § 4).
\item
  If \(n \leqslant c\) , let \(\mathbf{F}_{\mathbf{P},n}^{\mathrm{c}}\) be the intersection of \(\mathbf{F}_{\mathbf{P}}^{\mathrm{c}}\) with \(\Gamma_{n}/\Gamma_{c}\) ; if \(n \geqslant c\) , we write \(\mathbf{F}_{\mathbf{P},n}^{\mathrm{c}} = \{e\}\) . Show that \(\mathbf{F}_{\mathbf{P},n}^{\mathrm{c}}\) is the P-saturation of \(\mathbf{C}^{\mathrm{c}}\mathbf{F}^{\mathrm{c}}\) in \(\Gamma/\Gamma_{c}\) . The filtration \((\mathbf{F}_{\mathbf{P},n}^{\mathrm{c}})\) is an integral central filtration of \(\mathbf{F}_{\mathbf{P},n}^{\mathrm{c}}\) ; let \(\text{gr}(\mathbf{F}_{\mathbf{P}}^{\mathrm{c}})\) be the graded module associated with this filtration. Show that, if \(n < c\) , the image of \(\text{gr}_{\mathbf{n}}(\mathbf{F}_{\mathbf{P}}^{\mathrm{c}})\) in \(\text{gr}_{\mathbf{n}}(\Gamma) = \mathbf{A}_{\mathbf{Z}_{\mathbf{P}}}^{n}(\mathbf{X})\) is \(\mathbf{S}_{\mathbf{P}}^{-1}. \mathbf{L}_{\mathbf{Z}}^{n}(\mathbf{X}) = \mathbf{L}_{\mathbf{Z}_{\mathbf{P}}}^{n}(\mathbf{X})\) . Deduce that the Lie algebra \(\text{gr}(\mathbf{F}_{\mathbf{P},n}^{\mathrm{c}})\) is generated by its elements of degree 1 and hence that \(\mathbf{F}_{\mathbf{P},n} = \mathbf{C}^{n}(\mathbf{F}_{\mathbf{P}}^{\mathrm{c}})\) for all \(n\) ( \(\S 4\) , Exercise 1). Show that the group \(\mathbf{F}_{\mathbf{P}}^{\mathrm{c}}\) is generated by the \(x^{1/s}\) for \(x \in \mathbf{X}\) , \(s \in \mathbb{S}_{\mathbf{P}}\) (note that the images of these elements in \(\text{gr}_1(\mathbf{F}_{\mathbf{P}}^{\mathrm{c}}) \simeq \mathbf{Z}_{\mathbf{P}}^{(X)}\) generate the group \(\text{gr}_1(\mathbf{F}_{\mathbf{P}}^{\mathrm{c}})\) and apply Corollary 3 to Proposition 8 of \(Algebra\) , Chapter I, § 6, no. 3).
\item
  Let \(\mathrm{H}\) be a Hall set relative to \(\mathbf{X}\) (§2, no. 10) and let \(\mathrm{H}(c)\) be the subset of \(\mathrm{H}\) consisting of the elements of length \(< c\) . For \(m \in \mathrm{H}(c)\) , let \(\phi_c(m)\) denote the image in \(\mathbf{F}^c\) of the basic commutator \(\phi(m)\) defined by \(m\) (cf.~no. 4, Remark). Show that, for all \(w \in \mathbf{F}_P^c\) , there exists a unique element \(\alpha \in \mathbf{Z}_{\mathbf{P}}^{\mathbf{H}(c)}\) such that
\end{enumerate}

\[
w = \prod_ {m \in H (c)} \phi_ {c} (m) ^ {\alpha (m)}.
\]

(Use the determination of gr(F \(_{P}^{c}\) ) obtained above.)

\begin{enumerate}
\def\labelenumi{\arabic{enumi}.}
\setcounter{enumi}{6}
\tightlist
\item
  We preserve the notation of Exercise 6.
\end{enumerate}

\begin{enumerate}
\def\labelenumi{(\alph{enumi})}
\item
  Let G be a nilpotent group of class \(<c\) and let \((i, \overline{G})\) be a P-envelope of G. Let \(\rho: \mathrm{F}^{c} \to \mathrm{G}\) be a surjective homomorphism (such a homomorphism exists if X is suitably chosen) and let \(\bar{\rho}\) be the corresponding homomorphism of \(\mathrm{F}_{\mathbf{P}}^{c} \in \overline{G}\) (§ 4, Exercise 15). The homomorphism \(\bar{\rho}\) is surjective and maps \(\mathrm{C}^{n}\mathrm{F}_{\mathbf{P}}^{c} = \mathrm{F}_{\mathbf{P}, n}^{c}\) onto \(\mathrm{C}^{n}\overline{G}\) . Show that \(\mathrm{C}^{n}\overline{G}\) is the P-envelope of \(i(\mathrm{C}^{n}\mathrm{G})\) in \(\overline{G}\) and that the Lie algebra \(\mathrm{gr}(\overline{G})\) is identified with \(\mathbf{Z}_{\mathbf{P}} \otimes \mathrm{gr}(\mathrm{G})\) .
\item
  Deduce that G is P-divisible and P-torsion-free if and only if the \(C^{n}G/C^{n+1}G\) are.
\end{enumerate}

¶ 8. Suppose that K = Z and that X is finite. For every integer \(k \geqslant 0\) , let \((e_{k,\alpha})\) denote the basis of \(A_k(X)\) .

\begin{enumerate}
\def\labelenumi{(\alph{enumi})}
\tightlist
\item
  Let \((w_{n})_{n\in \mathbf{N}}\) be a sequence of elements of F. Let \(w_{k,\alpha}(n)\) be the coefficient of \(e_{k,\alpha}\) in the term of \(g(w_{n})\in \hat{\mathrm{A}} (\mathrm{X})\) of degree \(k\) . Then:
\end{enumerate}

\[
g (w _ {n}) = \sum_ {k, \alpha} w _ {k, \alpha} (n) e _ {k, \alpha} \quad \text { for   all } n \in \mathbf {Z}.
\]

The sequence \((w_{n})\) is called typical if, for every ordered pair \((k,\alpha)\) the function \(w_{k,\alpha}\colon n\mapsto w_{k,\alpha}(n)\) is a binomial polynomial of degree \(\leqslant k\) , cf.~Exercise 4. This condition is independent of the choice of the bases \((e_{k,\alpha})\) . Show that it is equivalent to the existence of a sequence \((a_{k})_{k\in \mathbf{N}}\) of elements of \(\hat{\mathbf{A}}_{\mathbf{q}}(\mathbf{X})\) such that \(\omega (a_k)\geqslant k\) for all \(k\) and

\[
g (w _ {n}) = \sum_ {k = 0} ^ {\infty} n ^ {k} a _ {k} \quad \text { for   all } n \in \mathbf {Z}.
\]

\begin{enumerate}
\def\labelenumi{(\alph{enumi})}
\setcounter{enumi}{1}
\item
  Show that, if \((w_{n})\) and \((w_{n}^{\prime})\) are two typical sequences, so are \((w_{n}^{-1})\) and \((w_{n}w_{n}^{\prime})\) .
\item
  Let \(w\) be an element of \(\mathbf{C}^k\mathbf{F}\) for \(k \geqslant 1\) and let \(f\) be a binomial polynomial of degree \(\leqslant k\) . Show that \((w^{f(n)})\) is a typical sequence. In particular if \(z \in \mathbf{F}\) , the sequence \((z^n)\) of powers of \(z\) is typical.
\item
  Let \((w_{n})\) be a sequence of elements of F. Show that \((w_{n})\) is typical if and only if there exist \(y_{0}, y_{1}, \ldots, y_{n}, \ldots\) in F, where \(y_{i} \in \mathbf{C}^{i}\mathbf{F}\) for all \(i \geqslant 1\) and
\end{enumerate}

\[
w _ {n} \equiv y _ {0} y _ {1} ^ {n} \dots y _ {k} ^ {\binom {n} {k}} \mod \mathrm{C} ^ {k + 1} \mathrm{F}
\]

for all \(n \in \mathbf{Z}\) and all \(k \geqslant 1\) . (Determine the \(y_{i}\) from \((w_{n})\) by setting successively \(n = 0, 1, \ldots\) . For example \(w_{0} = y_{0}, w_{1} = y_{0}y_{1}, \ldots\) )

\begin{enumerate}
\def\labelenumi{(\alph{enumi})}
\setcounter{enumi}{4}
\tightlist
\item
  Let H be a Hall set relative to X and let \((w_{n})\) be a sequence of elements of F. Let \((\alpha(m,n))_{m\in\mathbb{H},n\in\mathbf{Z}}\) be the family of integers such that
\end{enumerate}

\[
w _ {n} \equiv \prod_ {m \in \mathrm{H}} \phi (m) ^ {\alpha (m, n)} \mod \mathrm{C} ^ {c} \mathrm{F}
\]

for all \(n \in \mathbf{Z}\) and all \(c \geqslant 1\) (cf.~no. 4, Remark). Show that \((w_n)\) is typical if and only if, for all \(m \in \mathbf{H}\) , the function \(n \mapsto \alpha(m, n)\) is a binomial polynomial of degree \(\leqslant l(m)\) , where \(l(m)\) is the length of \(m\) .†

\begin{enumerate}
\def\labelenumi{(\alph{enumi})}
\setcounter{enumi}{5}
\tightlist
\item
  A sequence \((w_{n})\) is called 1-typical if the corresponding functions \(w_{k,\alpha}\) are binomial polynomials of degree \(\leqslant k - 1\) . Show that, if \((w_{n})\) is 1-typical, there exist \(y_{i} \in \mathbf{C}^{i + 1}\mathbf{F}\) for \(i = 0, 1, \ldots\) , such that
\end{enumerate}

\[
w _ {n} \equiv y _ {0} y _ {1} ^ {n} \dots y _ {k} ^ {\binom {n} {k}} \mod \mathrm{C} ^ {k + 1} \mathrm{F}
\]

for all \(n \in Z\) and all \(k \geqslant 1\) .

¶ 9. Let X be the set with two elements \{x, y\}.

\begin{enumerate}
\def\labelenumi{(\alph{enumi})}
\tightlist
\item
  Show that there exists a sequence \(w_{2}, w_{3}, \ldots\) of elements of \(\mathbf{F}\) such that \(w_{i} \in \mathbb{C}^{i}\mathbb{F}\) for all \(i\) and
\end{enumerate}

\[
(x y) ^ {n} \equiv x ^ {n} y ^ {n} w _ {2} ^ {\binom {n} {2}} \dots w _ {i} ^ {\binom {n} {i}} \mod \mathrm{C} ^ {i + 1} \mathrm{F}
\]

for all \(n \in \mathbf{Z}\) and all \(i \geqslant 1\) . (Apply Exercise 8 ( \(d\) ) to the sequence of \(x^{-n}(xy)^n\) .) Show that \(w_2 = y^{-1}(y^{-1}, x^{-1})y \equiv (x, y)^{-1} \mod \mathrm{C}^3\mathrm{F}\) .

\begin{enumerate}
\def\labelenumi{(\alph{enumi})}
\setcounter{enumi}{1}
\tightlist
\item
  Let G be a nilpotent group of class \(\leqslant c\) and let x, y be in G. Deduce from the above the following formula (``Hall's formula'');
\end{enumerate}

\[
(x y) ^ {n} = x ^ {n} y ^ {n} \prod_ {i = 2} ^ {c} w _ {i} (x, y) ^ {\binom {n} {i}}.
\]

\begin{enumerate}
\def\labelenumi{(\alph{enumi})}
\setcounter{enumi}{2}
\tightlist
\item
  Let \(p\) be a prime number. Show that there exists \(u_{i} \in \mathbf{C}^{i}\mathbf{F}\) ( \(2 \leqslant i \leqslant p - 1\) ) and \(w \in \mathbf{C}^{p}\mathbf{F}\) such that
\end{enumerate}

\[
(x y) ^ {p} = x ^ {p} y ^ {p} u _ {2} ^ {p} \dots u _ {p - 1} ^ {p} w.
\]

(Use (a), noting that the \(\binom{p}{i}\) is divisible by \(p\) if \(1 < i < p\) .) Let \(\bar{w}\) be the image of \(w\) in \(\mathrm{gr}_p(\mathsf{F}) = \mathrm{L}_{\mathbf{Z}}^p (\mathbf{X})\) and let \(\bar{w}_p\) be the image of \(\bar{w}\) in \(\mathrm{L}_{\mathbf{F}_p}^p (\mathbf{X})\) . Show that \(\bar{w}_p = \Lambda_p(x,y)\) , cf.~Chapter I, § 1, Exercise 19. (Use the extension \(g\) of F to \(\tilde{\mathbf{A}}(\mathbf{X})\) and compare the terms of degree \(p\) in \(g((xy)^p)\) and \(g(x^p y^p u_2^p \ldots u_{p-1}^p w)\) ). The first is equal to \((x + y)^p\) and the second is congruent mod. \(p\) to \(x^p + y^p + \bar{w}_p\) . Hence the result.)

Show that there exist \(a \in \mathbf{C}^2\mathbf{F}\) and \(w' \in \mathbf{C}^p\mathbf{F}\) such that

\[
(x y a) ^ {p} = x ^ {p} y ^ {p} w ^ {\prime}.
\]

(Use the above formula for \((xy)^p\) .) Deduce that in a nilpotent group of class \(p\) the \(p\) -th powers form a subgroup.

\begin{enumerate}
\def\labelenumi{(\alph{enumi})}
\setcounter{enumi}{3}
\tightlist
\item
  Show that there exists a sequence \(v_{3}, v_{4}, \ldots\) of elements of \(\mathbf{F}\) , where \(v_{i} \in \mathbf{C}^{i}\mathbf{F}\) for all \(i\) and
\end{enumerate}

\[
(x ^ {n}, y) \equiv (x, y) ^ {n} v _ {3} ^ {\binom {n} {2}} \dots v _ {i + 1} ^ {\binom {n} {i}} \mod \mathrm{C} ^ {i + 2} \mathrm{F}
\]

for all \(n \in \mathbf{N}\) and all \(i \geqslant 1\) . (Apply Exercise 8 (f) to the sequence of \((x^n, y)\) .)

Deduce that, if \(p\) is a prime number, there exists \(t_i \in \mathrm{C}^i\mathrm{F}\) for \(3 \leqslant i \leqslant p\) and \(z \in \mathrm{C}^{p+1}\mathrm{F}\) such that

\[
(x ^ {p}, y) = (x, y) ^ {p} t _ {3} ^ {p} \dots t _ {p} ^ {p} z.
\]

Show that the image \(\bar{z}_p\) of \(z\) in \(\mathrm{L}_{\mathbf{F}_p}^{p + 1}(\mathbf{X})\) is (ad \(x)^p (y)\) . (Same method as for \((c).)\)

¶ 10. Let p be a prime number and let G be a group with a real central filtration (G \(_{\alpha}\) ). Suppose that the relation x ∈ G \(_{\alpha}\) implies x \(^{p}\) ∈ G \(_{p\alpha}\) , in which case the filtration (G \(_{\alpha}\) ) is called restricted.

\begin{enumerate}
\def\labelenumi{(\alph{enumi})}
\item
  Show that the Lie algebra \(\operatorname{gr}(\mathbf{G})\) associated with \((\mathbf{G}_{\alpha})\) is such that \(p.\operatorname{gr}(\mathbf{G}) = 0\) and hence can be given an algebra structure over \(\mathbf{F}_p\) .
\item
  Let \(\xi \in \mathrm{gr}_{\alpha}(\mathbf{G})\) and let \(x\) be a representative of \(\xi\) in \(\mathbf{G}_{\alpha}\) . Show that the image of \(x^{p}\) in \(\mathrm{gr}_{p\alpha}(\mathbf{G})\) does not depend on the choice of \(x\) (use Exercise 9 (c)). If this image is denoted by \(\xi^{[p]}\) , prove that \(\xi \mapsto \xi^{[p]}\) is a linear mapping of \(\mathrm{gr}_{\alpha}(\mathbf{G})\) into \(\mathrm{gr}_{p\alpha}(\mathbf{G}')\) and that \((\xi + \xi')^{[p]} = \xi^{[p]} + \xi'^{[p]} + \Delta_p(\xi, \xi')\) . (Same method.)
\end{enumerate}

Show that, if \(\xi \in \mathrm{gr}_{\alpha}(\mathbf{G})\) and \(\eta \in \mathrm{gr}_{\beta}(\mathbf{G})\) , then

\[
[ \xi^ {[ p ]}, \eta ] = (\operatorname{ad} \xi) ^ {p} (\eta).
\]

(Use Exercise 9 (d).)

\begin{enumerate}
\def\labelenumi{(\alph{enumi})}
\setcounter{enumi}{2}
\tightlist
\item
  Show that there exists one and only one \(p\) -mapping of \(\operatorname{gr}(\mathbf{G})\) into itself (Chapter I, § 1, Exercise 20) which extends the mappings \(\xi \mapsto \xi^{[p]}\) defined on the \(\operatorname{gr}_{\alpha}(\mathbf{G})\) . With this structure, \(\operatorname{gr}(\mathbf{G})\) is a Lie \(p\) -algebra, said to be associated with the restricted filtration \((\mathbf{G}_{\alpha})\) .
\end{enumerate}

¶ 11. (a) Let A be a filtered algebra satisfying the conditions of §4, no. 5 and let \(\Gamma = A^{*} \cap (1 + A_{0}^{+})\) . \(\Gamma\) is given the filtration induced by that on A (loc. cit., Proposition 2). Show that, if K is a field of characteristic \(p > 0\) , \((\Gamma_{\alpha})\) is a restricted filtration (Exercise 10) and the embedding of \(\mathrm{gr}(\Gamma)\) in \(\mathrm{gr(A)}\) defined in loc. cit., Proposition 3 is compatible with the Lie \(p\) -algebra structures on \(\mathrm{gr}(\Gamma)\) and \(\mathrm{gr(A)}\) .

\begin{enumerate}
\def\labelenumi{(\alph{enumi})}
\setcounter{enumi}{1}
\item
  Suppose that \(\mathbf{K} = \mathbf{F}_p\) . We apply the above to the filtered algebra \(\hat{\mathbf{A}}(\mathbf{K})\) , in which \(\mathrm{F} = \mathrm{F}(\mathrm{X})\) is embedded by means of the homomorphism \(g\) . The filtration \((\Gamma_n)\) of \(\Gamma\) induces a filtration \((\mathbb{F}_n)\) on \(\mathbf{F}\) , which is a restricted filtration. The Lie \(p\) -algebra \(\operatorname{gr}(\mathbf{F})\) is identified with a Lie \(p\) -subalgebra of \(\mathrm{A}(\mathbf{X})\) containing \(\mathbf{X}\) . Deduce that \(\operatorname{gr}(\mathbf{F})\) contains the free Lie \(p\) -algebra \(\mathrm{L}(\mathbf{X}, p)\) , cf.~§ 3, Exercise 4 (e).
\item
  Let \(\mathbf{H}\) be a Hall set relative to \(\mathbf{X}\) ; for every integer \(i\) , let \(\mathbf{H}_i\) be the set of elements of \(\mathbf{H}\) of length \(i\) . Let \(w \in \mathbf{F}\) and let \(\alpha_i \in \mathbf{Z}^{(\mathbf{H}_i)}\) be such that
\end{enumerate}

\[
w \equiv \prod_ {i = 1} ^ {n} \prod_ {m \in H_i} \phi (m) ^ {\alpha_ {i} (m)} \mod C ^ {n + 1} F
\]

for all \(n\) , cf.~no. 4. For all \(m \in \mathbf{H}\) , let \(l(m)\) denote the length of \(m\) and \(q(m, w)\) the largest power of \(p\) which divides \(\alpha_{1}(m)\) , where \(i = l(w)\) (if \(\alpha_{1}(m) = 0\) , we make the convention \(q(m, w) = +\infty\) ). Let \(\mathrm{N} = \inf_{m \in \mathbb{H}} (l(m).q(m, w))\) . Show that the \(\phi(m)^{\alpha_{1}(m)}\) belong to \(\mathrm{F}_{\mathrm{N}}\) and even to \(\mathrm{F}_{\mathrm{N+1}}\) if \(l(w).q(m, w) > \mathrm{N}\) . Show that, if \(l(w).q(m, w) = \mathrm{N}\) , the image of \(\phi(m)^{\alpha_{1}(m)}\) in \(\mathrm{gr}_{\mathrm{N}}(\hat{\mathrm{A}}(\mathbf{X})) = \mathrm{A}^{\mathrm{N}}(\mathbf{X})\) is a nonzero multiple of \(\overline{m}^{q(m, w)}\) , where \(\overline{m}\) is the element of \(\mathrm{L}(\mathbf{X})\) defined by \(m\) ( \(\S 2\) , no. 11). Now these elements are linearly independent ( \(\S 3\) , Exercise 4); deduce that the image of \(w\) in \(\mathrm{gr}_{\mathrm{N}}(\mathbf{F})\) is \(\neq 0\) and belongs to \(\mathrm{L}(\mathbf{X}, p)\) , whence the fact that \(\mathrm{gr}(\mathbf{F}) = \mathrm{L}(\mathbf{X}, p)\) .

\begin{enumerate}
\def\labelenumi{\arabic{enumi}.}
\setcounter{enumi}{11}
\tightlist
\item
  Let g be a Lie algebra over a field k of characteristic p \textgreater{} 0.
\end{enumerate}

\begin{enumerate}
\def\labelenumi{(\alph{enumi})}
\tightlist
\item
  Let \(h\) be an integer \(\leqslant p - 1\) . \(\mathfrak{g}\) is said to satisfy Engel's \(h\) -th condition if \((\operatorname{ad} x)^h(y) = 0\) for every ordered pair \((x, y)\) of elements of \(\mathfrak{g}\) . Show that this is equivalent to
\end{enumerate}

\[
\sum_ {\sigma \in \mathfrak {S} _ {h}} \operatorname{ad} (x _ {\sigma (1)}) \circ \dots \circ \operatorname{ad} (x _ {\sigma (h)}) = 0
\]

for \(x_1, \ldots, x_h\) in \(\mathfrak{g}\) (apply the formula \((\operatorname{ad} x)^h = 0\) to the linear combinations of the \(x_i\) ).


\begin{enumerate}
\def\labelenumi{(\alph{enumi})}
\setcounter{enumi}{1}
\item
  Let \(x, y\) be elements of \(\mathfrak{g}\) such that \(\Lambda_p(ax, by) = 0\) for all \(a, b\) in \(k\) . Show that, if \(\Lambda_p^{r,s}\) is the bihomogeneous component of \(\Lambda_p\) of bidegree \((r, s)\) , then \(\Lambda_p^{r,s}(x, y) = 0\) for every ordered pair \((r, s)\) . Deduce (taking \(r = p - 1\) , \(s = 1\) ) that \((\text{ad } x)^{p-1}(y) = 0\) . In particular, if \(\Lambda_p(x, y) = 0\) for \(x, y\) in \(\mathfrak{g}\) , the Lie algebra \(\mathfrak{g}\) satisfies Engel's \((p - 1)\) -th condition.
\item
  Let \(\mathfrak{c}\) be the centre of \(\mathfrak{g}\) . Suppose that \((\operatorname{ad} x)^p = 0\) for all \(x \in \mathfrak{g}\) . Show that \(\mathfrak{g} / \mathfrak{c}\) satisfies Engel's \((p - 1)\) -th condition. (Show that \(\Lambda_p(\operatorname{ad} x, \operatorname{ad} y) = 0\) for \(x, y\) in \(\mathfrak{g}\) and apply \((b)\) to the Lie algebra \(\operatorname{ad} \mathfrak{g} = \mathfrak{g} / c\) .)
\item
  Let G be a group such that \(x^{p}=1\) for all \(x\in G\) and let \((\mathrm{G}_{\alpha})\) be a real central filtration of G. The filtration \((\mathrm{G}_{\alpha})\) is restricted (Exercise 10) and \(\mathrm{gr}(\mathrm{G})\) is a Lie p-algebra whose p-mapping is zero. Deduce that \(\mathrm{gr}(\mathrm{G})\) satisfies Engel's \((p-1)\) -th condition.†
\end{enumerate}

\exercisesfor{6}

\begin{enumerate}
\def\labelenumi{\arabic{enumi}.}
\tightlist
\item
  Show that \(\exp (\mathrm{U}).\exp (\mathrm{V}).\exp (-\mathrm{U}) = \exp \left(\sum_{n = 0}^{\infty}\frac{1}{n!} (\operatorname {ad}\mathrm{U})^{n}(\mathrm{V})\right)\) . Deduce that
\end{enumerate}

\[
\mathrm{H} (\mathrm{U}, \mathrm{H} (\mathrm{V}, - \mathrm{U})) = \sum_ {n = 0} ^ {\infty} \frac {1}{n !} (\operatorname{ad} \mathrm{U}) ^ {n} (\mathrm{V}) = (\exp (\operatorname{ad} \mathrm{U})) (\mathrm{V}).
\]

\begin{enumerate}
\def\labelenumi{\arabic{enumi}.}
\setcounter{enumi}{1}
\item
  Show that \(\mathbf{H}(\mathbf{U},\mathbf{V}) = \sum_{n = 0}^{\infty}\frac{1}{n!} (\operatorname {ad}\mathbf{U})^{n}(\mathbf{H}(\mathbf{V},\mathbf{U}))\) . (Apply Exercise 1, noting that \(\mathbf{H}(\mathbf{U},\mathbf{H}(\mathbf{H}(\mathbf{V},\mathbf{U}), - \mathbf{U})) = \mathbf{H}(\mathbf{U},\mathbf{V}).\)
\item
  Let \(\mathbf{M}(\mathbf{U},\mathbf{V}) = \mathbf{H}(-\mathbf{U},\mathbf{U} + \mathbf{V})\) . Then
\end{enumerate}

\[
\exp (\mathrm{M} (\mathrm{U}, \mathrm{V})) = \exp (- \mathrm{U}). \exp (\mathrm{U} + \mathrm{V}).
\]

Let \(M_{1}(U, V)\) (resp. \(H_{1}(U, V)\) ) denote the sum of the bihomogeneous terms of the series M (resp. H) whose degree in V is 1.

\begin{enumerate}
\def\labelenumi{(\alph{enumi})}
\tightlist
\item
  Show that
\end{enumerate}

\[
\mathrm{M} _ {1} (\mathrm{U}, \mathrm{V}) = \sum_ {n = 0} ^ {\infty} \frac {(- 1) ^ {n}}{(n + 1) !} (\mathrm{adU}) ^ {n} (\mathrm{V}) = (f (\mathrm{adU})) (\mathrm{V})
\]

where \(f(\mathrm{T}) = (1 - e^{-\mathrm{T}}) / \mathrm{T}\) .

(Same argument as for the proof of Proposition 5. It can also be reduced to Proposition 5 by using the identity

\[
\exp (\mathrm{U}). \exp (\mathrm{M} (\mathrm{U}, \mathrm{V})). \exp (- \mathrm{U}) = \exp (\mathrm{H} (\mathrm{U} + \mathrm{V}, - \mathrm{U})),
\]

together with Exercise 1.)

\begin{enumerate}
\def\labelenumi{(\alph{enumi})}
\setcounter{enumi}{1}
\tightlist
\item
  Show that \(\mathrm{H}(\mathrm{U},\mathrm{M}(\mathrm{U},\mathrm{V})) = \mathrm{U} + \mathrm{V}\) . Deduce that
\end{enumerate}

\[
\mathrm{H} _ {1} (\mathrm{U}, \mathrm{M} _ {1} (\mathrm{U}, \mathrm{V})) = \mathrm{V}.
\]

\begin{enumerate}
\def\labelenumi{(\alph{enumi})}
\setcounter{enumi}{2}
\tightlist
\item
  Let \(g(\mathrm{T}) = 1 / f(\mathrm{T}) = 1 + \mathrm{T} / 2 + \sum_{n=1}^{\infty} \frac{1}{(2n)!} b_{2n}\mathrm{T}^{2n}\) , where the \(b_{2n}\) are
\end{enumerate}

the Bernoulli numbers (Functions of a Real Variable, Chapter VI, § 1, no. 4). Show that

\[
\mathrm{H} _ {1} (\mathrm{U}, \mathrm{V}) = (g (\operatorname{ad} \mathrm{U})) (\mathrm{V}) = \mathrm{V} + \frac {1}{2} [ \mathrm{U}, \mathrm{V} ] + \sum_ {n = 1} ^ {\infty} \frac {1}{(2 n) !} b _ {2 n} (\operatorname{ad} \mathrm{U}) ^ {2 n} (\mathrm{V}).
\]

Derive the first few terms of the expansion of \(H_{1}(U, V)\) :

\[
\begin{array}{r l} \mathrm{H} _ {1} (\mathrm{U}, \mathrm{V}) & = \mathrm{V} + \frac {1}{2} (\mathrm{adU}) (\mathrm{V}) + \frac {1}{12} (\mathrm{adU}) ^ {2} (\mathrm{V}) - \frac {1}{720} (\mathrm{adU}) ^ {4} (\mathrm{V}) \\ & + \frac {1}{30240} (\mathrm{adU}) ^ {6} (\mathrm{V}) - \dots \end{array}
\]

\begin{enumerate}
\def\labelenumi{(\alph{enumi})}
\setcounter{enumi}{3}
\tightlist
\item
  Show that the sum of the bihomogeneous terms of H(U, V) whose degree in U is 1 is the series:
\end{enumerate}

\[
\mathrm{U} + \frac {1}{2} [ \mathrm{U}, \mathrm{V} ] + \sum_ {n = 1} ^ {\infty} \frac {1}{(2 n) !} b _ {2 n} (\mathrm{adV}) ^ {2 n} (\mathrm{U}).
\]

(Apply (c) to H(V, U) and use Exercise 2 to pass from H(V, U) to H(U, V).) (e) Let W be another indeterminate and let X(U, V, W) be the sum of the trihomogeneous terms of H(U, V + W) whose degree in W is 1. Show that

\[
\mathrm{X} (\mathrm{U}, \mathrm{H} _ {1} (\mathrm{V}, \mathrm{W})) = \mathrm{H} _ {1} (\mathrm{H} (\mathrm{U}, \mathrm{V}), \mathrm{W}).
\]

(Use the identity \(\mathrm{H}(\mathrm{U},\mathrm{H}(\mathrm{V},\mathrm{W})) = \mathrm{H}(\mathrm{H}(\mathrm{U},\mathrm{V}),\mathrm{W}).\) ) Deduce that

\[
\mathrm{X} (\mathrm{U}, \mathrm{V}, \mathrm{W}) = (g (\operatorname{ad} \mathrm{H} (\mathrm{U}, \mathrm{V})) \circ f (\operatorname{ad} \mathrm{V})) (\mathrm{W}).
\]

¶ 4. Let X be a finite set and let \(\hat{\mathbf{L}} = \hat{\mathbf{L}}(\mathbf{X})\) . We give \(\hat{\mathbf{L}}\) the Hausdorff law of composition \((a, b) \mapsto a \uplus b\) , which we denote simply by \(a.b\) . If \(t \in \mathbf{K}\) , \(a \in \hat{\mathbf{L}}\) , we write \(a^t = t a\) .

\begin{enumerate}
\def\labelenumi{(\alph{enumi})}
\tightlist
\item
  Let \(\mathcal{H}\) be a Hall set relative to X. If \(m \in \mathcal{H}\) , let \(m_{\mathrm{L}}\) denote the corresponding element of the Lie algebra \(\mathrm{L}(\mathrm{X})\) (§2, no. 11) and \(m_{\mathcal{H}}\) the basic commutator of the group \(\hat{\mathrm{L}}\) defined by \(m\) (§5, no. 4, Remark). If \(m\) is of length \(l\) , show that \(m_{\mathcal{H}} = m_{\mathrm{L}} + \mu\) , where \(\mu\) is of filtration \(\geqslant l + 1\) (argue by induction on \(l\) ). Deduce that every element \(w\) of the group \(\hat{\mathrm{L}}\) can be written uniquely in the form
\end{enumerate}

\[
w = \prod_ {m \in \mathcal {H}} m _ {\mathcal {H}} ^ {\alpha (m)},
\]


where \(\alpha(m) \in \mathbf{K}\) and the product is convergent in the topological group \(\hat{\mathbf{L}}\) .

\begin{enumerate}
\def\labelenumi{(\alph{enumi})}
\setcounter{enumi}{1}
\item
  Let \(\mathbf{P}\) be the set of prime numbers and let \(c\) be an integer \(\geqslant 1\) . Let \(\mathbf{K} = \mathbf{Q}\) . Show that the group \(\hat{\mathbf{L}} / \hat{\mathbf{L}}_c\) is identified with the \(\mathbf{P}\) -envelope \(\mathrm{F}_{\mathbf{P}}^{c}\) of the group \(\mathrm{F}^c = \mathrm{F}(\mathrm{X}) / \mathrm{C}^c\mathrm{F}(\mathrm{X})\) , cf.~§ 5, Exercise 6.
\item
  Take X to be a set of two elements \(\{U, V\}\) . Show the existence of two families \(\alpha(m)\) , \(\beta(m)\) of rational numbers such that
\end{enumerate}

\[
\mathrm{U} + \mathrm{V} = \prod_ {m \in \mathcal {H}} m _ {\mathcal {H}} ^ {\alpha (m)}
\]

\[
[ \mathrm{U}, \mathrm{V} ] = \prod_ {m \in \mathcal {H}} m _ {\mathcal {H}} ^ {\beta (m)}.
\]

For example

\[
\begin{array}{c} \mathrm{U} + \mathrm{V} = \mathrm{U}. \mathrm{V}. (\mathrm{U}, \mathrm{V}) ^ {- \frac {1}{2}} \dots \\ {} [ \mathrm{U}, \mathrm{V} ] = (\mathrm{U}, \mathrm{V}) \dots \end{array}
\]

\begin{enumerate}
\def\labelenumi{(\alph{enumi})}
\setcounter{enumi}{3}
\tightlist
\item
  Let g be a nilpotent Lie Q-algebra with the Hausdorff group law. Let u, v be in g. Show that
\end{enumerate}

\[
u + v = \prod_ {m \in \mathcal {H}} m _ {\mathcal {H}} (u, v) ^ {\alpha (m)}
\]

\[
[ u, v ] = \prod_ {m \in \mathcal {H}} m _ {\mathcal {H}} (u, v) ^ {\beta (m)},
\]

where \(\alpha\) and \(\beta\) are the families of rational numbers defined above (``inversion of the Hausdorff formula''). (Use the continuous homomorphism \(\phi\colon\hat{L}\to\mathfrak{g}\) such that \(\phi(U)=u\) , \(\phi(V)=v\) and note that it is also a homomorphism for the Hausdorff group law.)

\begin{enumerate}
\def\labelenumi{(\alph{enumi})}
\setcounter{enumi}{4}
\tightlist
\item
  Let G be a torsion-free divisible nilpotent group (i.e.~P-torsion-free and P-divisible). If \(u \in G\) , \(t \in Q\) , \(u^{t}\) is defined ( \(\S 4\) , Exercise 16). If \(u \in G\) , \(v \in G\) , we define \(u + v\) and {[}u, v{]} by the formulae of (d) above. Show that G thus has a nilpotent Lie Q-algebra structure and that the corresponding Hausdorff law is the given group law on G (verify these assertions when G is a group \(F_{P}^{c}\) , cf.~(b)), and pass from this to the general case using homomorphisms of \(F_{P}^{r}\) into G).
\end{enumerate}

Let \(f\) be a mapping of \(G\) into a torsion-free divisible nilpotent group \(G'\) . Show that \(f\) is a (group) homomorphism if and only if it is a Lie algebra homomorphism. (The Hausdorff law therefore defines an isomorphism of the ``category'' of nilpotent Lie \(\mathbf{Q}\) -algebras onto that of torsion-free divisible nilpotent groups.)

\begin{enumerate}
\def\labelenumi{\arabic{enumi}.}
\setcounter{enumi}{4}
\tightlist
\item
  Let \(\mathfrak{g}\) be a nilpotent Lie \(\mathbf{Q}\) -algebra which we give the Hausdorff group law, denoted by \((x,y)\mapsto x.y\) .
\end{enumerate}

\begin{enumerate}
\def\labelenumi{(\alph{enumi})}
\tightlist
\item
  Show that, if \(x \in \mathfrak{g}\) , \(y \in \mathfrak{g}\) , then
\end{enumerate}

\[
x. y. x ^ {- 1} = \sum_ {n = 0} ^ {\infty} \frac {1}{n !} (\operatorname{ad} x) ^ {n} (y).
\]

(Use Exercise 2.)

\begin{enumerate}
\def\labelenumi{(\alph{enumi})}
\setcounter{enumi}{1}
\item
  Let \(\mathfrak{h}\) be a subset of \(\mathfrak{g}\) . Show that \(\mathfrak{h}\) is a Lie subalgebra (resp. an ideal) of \(\mathfrak{g}\) if and only if it is a saturated subgroup (resp. a saturated normal subgroup) of the group \(\mathfrak{g}\) . (Use the formulae of Exercise 4 (d) to pass from the group law to the Lie algebra law.)
\item
  Let \(\mathfrak{h}\) be a subgroup of the group \(\mathfrak{g}\) . Show that the \(\mathbf{P}\) -saturation of \(\mathfrak{h}\) is \(\mathbf{Q}.\mathfrak{h}\) .
\item
  Let \(\mathfrak{h}\) be a Lie subalgebra of \(\mathfrak{g}\) . Show that the centralizer (resp. normalizer) of \(\mathfrak{h}\) in the group \(\mathfrak{g}\) is the set of \(x \in \mathfrak{g}\) such that \((\operatorname{ad} x)(\mathfrak{h}) = 0\) (resp. \((\operatorname{ad} x)(\mathfrak{h}) \subset \mathfrak{h}\) ).
\item
  Show that the lower central series of the Lie algebra g coincides with that of the group g and that the associated graded Lie algebra gr(g) is the same from the ``group'' point of view as from the ``Lie algebra'' point of view.
\end{enumerate}

¶ 6. Let G be a finitely generated torsion-free nilpotent group. Show that, if n is sufficiently large, G can be embedded in the lower strict triangular group of order n over Z. (Let \((i, \overline{G})\) be a P-envelope of G, cf.~§ 5, Exercise 6; with its canonical Lie algebra structure (Exercise 4), \(\overline{G}\) is a finite-dimensional nilpotent Lie Q-algebra. Apply Ado's Theorem (Chapter 1, § 7) to \(\overline{G}\) and deduce an embedding of \(\overline{G}\) in a strict triangular group over Q. Pass from Q to Z by conjugating by a suitable integral diagonal matrix.)

\exercisesfor{7}

\begin{enumerate}
\def\labelenumi{\arabic{enumi}.}
\tightlist
\item
  Let \(\mathfrak{g}\) be a complete normed Lie algebra over \(\mathbf{K}\) such that
\end{enumerate}

\[
\| [ x, y ] \| \leqslant \| x \| \| y \|
\]

for \(x, y\) in \(\mathfrak{g}\) . Let \(\Theta\) denote the set of \(x \in \mathfrak{g}\) such that \(\| x \| < \frac{1}{3} \log \frac{3}{2}\) . For \(x, y\) in \(\Theta\) , \(h(x, y) \in \Theta\) , cf.~no. 2.

\begin{enumerate}
\def\labelenumi{(\alph{enumi})}
\item
  We write \(f(\mathrm{T}) = (1 - \varepsilon^{-\mathrm{T}}) / \mathrm{T}\) and \(g(\mathrm{T}) = 1 / f(\mathrm{T})\) . The series \(f\) converges throughout the complex plane and the series \(g\) on the open disc of radius \(2\pi\) (cf.~Functions of a Real Variable, Chapter VI, §2, no. 3). Deduce that, if \(z \in \Theta\) , \(f(\mathrm{ad}z)\) and \(g(\mathrm{ad}z)\) are defined and that they are elements of \(\mathcal{L}(\mathfrak{g};\mathfrak{g})\) and inverses of one another.
\item
  Let \(\mathrm{D}_2h(x,y)\) be the second partial derivative of \(h\) at a point \((x,y)\) of \(\Theta \times \Theta\) (Differentiable and Analytic Manifolds, R, 1.6.2). It is an element of \(\mathcal{L}(\mathfrak{g};\mathfrak{g})\) . Show that
\end{enumerate}

\[
\mathrm{D} _ {2} h (x, y) = g (\operatorname{ad} h (x, y)) \circ f (\operatorname{ad} y).
\]

(Use Exercise 3 (e) of §6 to show this formula when \(x\) and \(y\) are sufficiently close to zero and pass from this to the general case by analytic continuation.) Show that the formula

\[
f (\operatorname{ad} h (x, y)) \circ \mathrm{D} _ {2} h (x, y) = f (\operatorname{ad} y)
\]

is valid at every point of the domain of absolute convergence of the formal power series \(\tilde{H}\) (cf.~Proposition 1).

\exercisesfor{8}

We suppose that the residual characteristic \(p\) of the field \(K\) is \(>0\) . We denote by \(\mathfrak{o}_{K}\) the ring of the valuation \(v\) of \(K\) .

\begin{enumerate}
\def\labelenumi{\arabic{enumi}.}
\tightlist
\item
  \begin{enumerate}
  \def\labelenumii{(\alph{enumii})}
  \tightlist
  \item
    An element \(\pi\) of \(K\) is called admissible if \(v(\pi) = \theta\) and \(\pi^{p-1} / p \equiv -1 \pmod{\pi o_K}\) .
  \end{enumerate}
\end{enumerate}

(Note that \(\pi^{p-1}/p\in\mathfrak{o}_{K}\) since \(v(\pi)=\theta.\) )

Show that there exist fields K satisfying the conditions of the paragraph and containing an admissible element (adjoin to \(Q_{p}\) a \((p - 1)\) -th root of \(-p\) ).

\begin{enumerate}
\def\labelenumi{(\alph{enumi})}
\setcounter{enumi}{1}
\tightlist
\item
  Let \(\pi\) be an admissible element of K. Consider the following formal power series with coefficients in K:
\end{enumerate}

\[
e _ {\pi} (\mathrm{U}) = \frac {1}{\pi} e (\pi \mathrm{U}) = \sum_ {n = 1} ^ {\infty} \pi^ {n - 1} \mathrm{U} ^ {n} / n!,
\]

\[
l _ {\pi} (\mathrm{U}) = \frac {1}{\pi} l (\pi \mathrm{U}) = \sum_ {n = 1} ^ {\infty} (- \pi) ^ {n - 1} \mathrm{U} ^ {n} / n.
\]

Show that these series are inverses of one another and that their coefficients belong to \(\mathfrak{o}_{K}\) (use Lemma 2).

\begin{enumerate}
\def\labelenumi{(\alph{enumi})}
\setcounter{enumi}{2}
\tightlist
\item
  Let \(\{\mathrm{U},\mathrm{V}\}\) be a set with two elements and let \(\mathrm{H}(\mathrm{U},\mathrm{V})\) be the Hausdorff series. We write
\end{enumerate}

\[
\mathrm{H} _ {\pi} (\mathrm{U}, \mathrm{V}) = \frac {1}{\pi} \mathrm{H} (\pi \mathrm{U}, \pi \mathrm{V}).
\]

Then \(H_{\pi}(U, V) \in \hat{L}_{K}(\{U, V\})\) .

Show that

\[
\mathrm{H} _ {\pi} (\mathrm{U}, \mathrm{V}) = l _ {\pi} (e _ {\pi} (\mathrm{U}) + e _ {\pi} (\mathrm{V}) + \pi e _ {\pi} (\mathrm{U}). e _ {\pi} (\mathrm{V})).
\]

Deduce that \(H_{\pi} \in \hat{L}_{\mathfrak{o}_{K}}(\{\mathrm{U}, \mathrm{V}\})\) , i.e.~that the coefficients of \(H_{\pi}\) belong to \(o_{K}\) ; hence another proof of Proposition 1.

\begin{enumerate}
\def\labelenumi{(\alph{enumi})}
\setcounter{enumi}{3}
\tightlist
\item
  Let \(\tilde{e}_{\pi}, \tilde{l}_{\pi}\) and \(\tilde{\mathbf{H}}_{\pi}\) denote the series obtained by reducing \(e_{\pi}, l_{\pi}\) and \(\mathbf{H}_{\pi}\) modulo \(\pi \mathfrak{o}_{K}\) ; their coefficients belong to the ring \(\mathfrak{o}_{K} / \pi \mathfrak{o}_{K}\) .
\end{enumerate}

Show that \(\tilde{l}_{\pi}(\mathrm{U}) = \mathrm{U} - \mathrm{U}^{p}\) (note that \(v_{p}(n) < (n - 1)\theta\) if \(n \neq 1, p\) ). Deduce that

\[
\tilde {e} _ {\pi} (\mathrm{U}) = \mathrm{U} + \mathrm{U} ^ {p} + \dots + \mathrm{U} ^ {p ^ {n}} + \dots
\]

and that

\[
\tilde {\mathrm{H}} _ {\pi} (\mathrm{U}, \mathrm{V}) = \tilde {e} _ {\pi} (\mathrm{U}) + \tilde {e} _ {\pi} (\mathrm{V}) - (\tilde {e} _ {\pi} (\mathrm{U}) + \tilde {e} _ {\pi} (\mathrm{V})) ^ {p},
\]

or also

\[
\tilde {\mathrm{H}} _ {\pi} (\mathrm{U}, \mathrm{V}) = \mathrm{U} + \mathrm{V} - \Lambda_ {p} \left(\sum_ {n = 0} ^ {\infty} \mathrm{U} ^ {p ^ {n}}, \sum_ {n = 0} ^ {\infty} \mathrm{V} ^ {p ^ {n}}\right),
\]

where \(\Lambda_{p}\) is defined by the formula

\[
\Lambda_ {p} (\mathrm{U}, \mathrm{V}) = (\mathrm{U} + \mathrm{V}) ^ {p} - \mathrm{U} ^ {p} - \mathrm{V} ^ {p}, \quad \text { cf.   Chapter   I,   } \S 1, \text {   Exercise   19. }
\]

In particular, \(\tilde{\mathrm{H}}_{\pi}(\mathrm{U},\mathrm{V})\) belongs to \(\hat{\mathrm{L}}_{\mathbf{F}_{p}}(\{\mathrm{U},\mathrm{V}\})\) and its homogeneous component of degree \(p\) is \(-\Lambda_p(\mathrm{U},\mathrm{V})\) . Then

\[
\begin{array}{c} \tilde {\mathrm{H}} _ {\pi} (\mathrm{U}, \mathrm{V}) = \tilde {\mathrm{H}} _ {\pi} (\mathrm{V}, \mathrm{U}) \\ \tilde {\mathrm{H}} _ {\pi} (\mathrm{U}, \tilde {\mathrm{H}} _ {\pi} (\mathrm{V}, \mathrm{W})) = \tilde {\mathrm{H}} _ {\pi} (\tilde {\mathrm{H}} _ {\pi} (\mathrm{U}, \mathrm{V}), \mathrm{W}) \quad \text { in } \hat {\mathrm{L}} _ {\mathbf {F} _ {p}} (\{\mathrm{U}, \mathrm{V}, \mathrm{W} \}), \end{array}
\]

where W is a third indeterminate.

\begin{enumerate}
\def\labelenumi{(\alph{enumi})}
\setcounter{enumi}{4}
\tightlist
\item
  Let C(U, V) = H(-U, H(-V, H(U, V))) (``Hausdorff commutator''). Then exp(C(U, V)) = exp(-U) exp(-V) exp(U) exp(V). We write
\end{enumerate}

\[
\mathrm{C} _ {\pi} (\mathrm{U}, \mathrm{V}) = \frac {1}{\pi} \mathrm{C} (\pi \mathrm{U}, \pi \mathrm{V}).
\]

Show that the coefficients of \(C_{\pi}(U,V)\) belong to \(\pi o_{K}\) (use the fact that \(\tilde{H}_{\pi}(U,V)=\tilde{H}_{\pi}(V,U)\) ).†

\begin{enumerate}
\def\labelenumi{\arabic{enumi}.}
\setcounter{enumi}{1}
\tightlist
\item
  We know (§ 6, Exercise 3) that if \(n \geqslant 2\) the component of H(U, V) of bidegree \((n, 1)\) is \(\frac{1}{n!} b_n(\mathrm{ad}\mathrm{U})^n (\mathrm{V})\) , where \(b_n\) is the \(n\) -th Bernoulli number.
\end{enumerate}

Deduce from this and Proposition 1 the inequality

\[
v _ {p} (b _ {n} / n!) \leqslant n / (p - 1).
\]

Recover this result by means of the Clausen-von Staudt Theorem (Functions of a Real Variable, Chapter VI, §2, Exercise 6) and show that there is equality if and only if \(S(n) = p - 1\) .

\begin{enumerate}
\def\labelenumi{\arabic{enumi}.}
\setcounter{enumi}{2}
\tightlist
\item
  Suppose that K contains a primitive \(p\) -th root of unity \(w\) . We write \(\pi = w - 1\) . Using the formula
\end{enumerate}

\[
w ^ {i} - 1 = \pi (1 + w + \dots + w ^ {i - 1})
\]

show that \(v(w^{i} - 1) = v(\pi)\) for \(1 \leqslant i \leqslant p - 1\) and that

\[
\frac {w ^ {i} - 1}{\pi} \equiv i \pmod {\pi}.
\]

Deduce, by means of the formula \(p = \prod_{i=1}^{p-1} (w^i - 1)\) , that \(\pi\) is admissible (Exercise 1).

¶ 4. Let c be an integer \(\geqslant1\) and let P be a set of prime numbers containing all prime numbers \(\leqslant c\) . Let \(Z_{P}=S_{P}^{-1}Z\) (§ 4, Exercise 16). Show that the terms of degree \(\leqslant c\) in the Hausdorff series H(U, V) belong to \(L_{Z_{P}}(\{U, V\})\) . Deduce that, if g is a nilpotent Lie \(Z_{P}\) -algebra of class \(\leqslant c\) , the law of composition \((u,v)\mapsto\mathrm{H}(u,v)\) makes g into a P-torsion-free P-divisible nilpotent group of class \(\leqslant c\) . Show that conversely every group with these properties can be obtained by this process. (Use ``the inversion of the Hausdorff formula'', cf.~§6, Exercise 4, and show that the exponents \(\alpha(m)\) , \(\beta(m)\) which appear there belong to \(Z_{P}\) when \(l(m)\leqslant c\) .)

In particular, every p-group of order \(p^{n}\) and class \textless p can be obtained by means of the Hausdorff law from a nilpotent Lie Z-algebra of class \textless p with \(p^{n}\) elements. (Take P to be the set of prime numbers distinct from p.)

\exercisesforappendix

\begin{enumerate}
\def\labelenumi{\arabic{enumi}.}
\tightlist
\item
  Let \(\Phi_n\) be the cyclotomic polynomial of index \(n\) (Algebra, Chapter V, § 11, no. 2). Using the formula
\end{enumerate}

\[
\mathbf {X} ^ {n} - 1 = \prod_ {d | n} \Phi_ {d} (\mathbf {X}),
\]

show that

\[
\Phi_ {n} (\mathrm{X}) = \prod_ {d | n} \left(\mathrm{X} ^ {n / d} - 1\right) ^ {\mu (d)}.
\]

(Apply the Möbius inversion formula to the multiplicative group of the field \(\mathbf{Q}(\mathbf{X})\) .)

\begin{enumerate}
\def\labelenumi{\arabic{enumi}.}
\setcounter{enumi}{1}
\tightlist
\item
  Let D be the total algebra of the monoid \(\mathbf{N}^*\) (Algebra, Chapter III, § 2, no. 10). If \(n \in \mathbf{N}^*\) , let \(n^\omega\) denote its image in D, so that \(1^\omega = 1\) and \((nm)^\omega = n^\omega m^\omega\) for \(n\) , \(m\) in \(\mathbf{N}^*\) . Every element \(f\) of D can be written uniquely as a series \(\sum_{i=1}^{\infty}\)
\end{enumerate}

\(\sum_{n=1}^{\infty} a_n n^\omega\) , where \(a_n \in \mathbf{K}\) .

\begin{enumerate}
\def\labelenumi{(\alph{enumi})}
\item
  Let \(f = \sum_{n=1}^{\infty} a_n n^\omega\) be an element of D. Show that \(f\) is invertible in D if and only if \(a_1\) is invertible in K. In particular, if K is a local ring, so is D.
\item
  We write \(\zeta = \sum_{n=1}^{\infty} n^{\omega}\) and \(\mu = \sum_{n=1}^{\infty} \mu(n)n^{\omega}\) . Show that \(\zeta\) and \(\mu\) are inverses of one another. Deduce that, if \(s = \sum_{n} s(n)n^{\omega}\) and \(t = \sum_{n} t(n)n^{\omega}\) are two elements of D, the relations \(s = \zeta\) . \(t\) and \(t = \mu\) . \(s\) are equivalent (variant of the Möbius inversion formula).
\item
  Let \(\mathbf{P}\) be the set of prime numbers. Show that the family of \((1 - p^{\omega})\) for \(p\in \mathbf{P}\) is multipliable in \(\mathbf{D}\) and that
\end{enumerate}

\[
\mu = \sum_ {p \in P} ^ {\infty} (1 - p ^ {\omega}) \quad \text { and } \quad \zeta = \prod_ {p \in P} \frac {1}{1 - p ^ {\omega}}.
\]

\chapter{Lie Groups}

Throughout the chapter, K denotes either the valued field R of real numbers, or the valued field C of complex numbers, or a non-discrete complete ultrametric commutative field. We assume that K is of characteristic 0 from § 4 onwards, that K = R or C in § 6, that K is ultrametric in § 7. Unless otherwise mentioned, all the manifolds, all the algebras and all the vector spaces considered are over K. Recall that, when we speak of a manifold of class \(C^{r}\), r ∈ N\_K, that is r = ω if K ≠ R and l ≤ r ≤ ω if K = R.

The conventions on norms, normable spaces and normed spaces are those of Differentiable and Analytic Manifolds, R.

Recall that a normable algebra over K is a (not necessarily associative) algebra A over K, with a topology T which possesses the following properties:

\begin{enumerate}
\def\labelenumi{(\arabic{enumi})}
\item
  \(\mathcal{T}\) can be defined by a norm;
\item
  the mapping \((x,y)\mapsto xy\) of \(\mathbf{A}\times \mathbf{A}\) into \(\mathbf{A}\) is continuous.
\end{enumerate}

The group of bicontinuous automorphisms of A is denoted by \(\mathrm{Aut(A)}\) . Every finite-dimensional algebra over \(K\) is a normable algebra with the canonical topology. A normed algebra over \(K\) is an algebra A over \(K\) with a norm such that \(\| xy\| \leqslant \| x\|\| y\|\) for all \(x,y\) in A; the algebra A with the topology defined by this norm is a normable algebra. If A is a normable algebra, there exists a norm on A defining its topology and making A into a normed algebra.

If G is a group, \(e_{G}\) , or simply e, denotes the identity element of G. For \(g \in G\) , \(\gamma(g)\) , \(\tilde{\delta}(g)\) and \(\operatorname{Int}(g)\) denote the mappings \(g' \mapsto gg'\) , \(g' \mapsto g'g^{-1}\) and \(g' \mapsto gg'g^{-1}\) of G into G. If f is a mapping of G into a set E, f denotes the mapping \(g \mapsto f(g^{-1})\) of G into E.

\section{LIE GROUPS}
